% Modelo de livro da série Ensino de Graduação da Editora UnB
% Autor: Leonardo Luiz e Castro

% Compile usando LuaLaTeX.

\documentclass{book}

% Chamadas e configurações de pacotes estão na pasta config:

% O comando abaixo pode ser necessário em documentos que usam versões antigas do pacote "fontspec". Se for começar um novo documento, deixe esse linha comentada e use caracteres Unicode.
% \usepackage[utf8]{luainputenc}

% Tamanho do papel e das margens: 

\usepackage[paperheight=230mm,paperwidth=160mm,inner=20mm,outer=14mm,top=24.5mm,bottom=21mm,pdftex]{geometry}

% Suporte a idiomas:

\usepackage[greek,brazilian]{babel}

% Para usar as variáveis \theauthor, \thetitle e \thedate (autor, título e data).

\usepackage{titling}

% Título de capítulo elegante:

\usepackage[Lenny]{fncychap}

% Pacote para estilo elegante:

\usepackage{fancyhdr}

% Definição de novas cores:

\usepackage{xcolor}

\definecolor{verde}{rgb}{0.2, 0.50, 0.25}
\definecolor{verde_UnB}{cmyk}{1,0,1,0.2}
\definecolor{cinza_UnB}{rgb}{0.6,0.6,0.6} % https://www.ginifab.com/feeds/pms/color_picker_from_image.php


% Para mudar cor de títulos (https://tex.stackexchange.com/a/75670/91816):

\usepackage{sectsty}

\chapterfont{\color{verde_UnB}}  % sets colour of chapters
\sectionfont{\color{verde_UnB}}  % sets colour of sections

% Pacote para definir fonte de qualquer tamanho:

\usepackage{fix-cm}

%% Pacote para correta separação silábica de palavras com hífen;
%% hífens devem ser escritos como \Hyphdash ou \-/ (se opção shortcuts estiver ativa)
%% no texto; por exemplo cana\-/de\-/açúcar seria separada como
%% cana-de- numa linha e -açúcar na linha seguinte
%% repetindo o hífen (um hífen é da separação e outro é da palavra):

\usepackage[shortcuts]{extdash}

%% Definindo espacamento:

\usepackage{setspace} 
\singlespacing

%% Enumerar itens:

\usepackage{enumitem}

%% Pacote para criar caixas:

\usepackage{pbox} %  que limitam a largura de textos em células,
%\usepackage{minibox} % que têm largura arbitrária.

%% Definição das dimensões do texto:

%\setlength{\textwidth}{16cm}
%\setlength{\textheight}{22cm}
%\setlength{\headheight}{1cm}
%\setlength{\footheight}{1cm}

\usepackage[title,titletoc]{appendix}

%% Passando títulos para o português:

\renewcommand{\chaptername}{Capítulo}
\renewcommand{\bibname}{Referências}
\renewcommand{\appendixname}{Apêndice}
\renewcommand{\indexname}{Índice Remissivo}
\renewcommand{\contentsname}{\bf\color{verde_UnB}Sumário}
\renewcommand{\tablename}{Tabela}
\renewcommand{\figurename}{\bf Fig.}
\renewcommand{\sin}{sen}
\def\listoffiguresname{Lista de Figuras}
\def\listoftablesname{Lista de Tabelas}

%% Definindo estilo de pagina:

\pagestyle{fancy}

% PACOTES DE TABELAS:

%% Espaçamento de tabelas ajustado:

\usepackage{booktabs}

%% Pacotes de linhas e colunas multiplas:

\usepackage{multirow}
\usepackage{multicol}

%% Pacotes de tabela e tabulação:

%%% Tabela com caixas centralizada:

\usepackage{array}
\newcolumntype{P}[1]{>{\centering\arraybackslash}p{#1}}

%%% Tabela em múltiplas páginas:

\usepackage{longtable}

%%% Tabela colorida:

\usepackage{tabu}
\usepackage{colortbl}

%%% Mudando cor das linhas de todas as tabelas:

\makeatletter
\renewenvironment{table}
     {\@float{table}\taburulecolor{verde_UnB}\arrayrulecolor{verde_UnB}}
     {\end@float}
\makeatother

%%% Tabela com linhas pontilhadas:

\usepackage{arydshln}

% PACOTES MATEMÁTICOS:

%% Pacote para não-itálico em ambiente de fórmulas:
%\usepackage[]{mathastext}

%% Pacotes matemáticos:

\usepackage{amsmath}
\usepackage{mathtools}

%%% Numera apenas equações usadas:

\mathtoolsset{showonlyrefs=true}

%%% Declarando barras para valor absoluto e norma:

\DeclarePairedDelimiter\abs{\lvert}{\rvert}%
\DeclarePairedDelimiter\norm{\lVert}{\rVert}%

%% Nome de equações ao lado com comando "eqname":

\newcommand{\eqname}[1]{\tag*{\llap{#1}}}

%% Fontes tipográficas:

\usepackage{fontspec}
%\usepackage{libertineotf}

%%% Seleção da fonte UnB Pro:

\setmainfont{UnB-Pro}[
  Path = fonts/UnB_Pro_v1.0/,
  Extension=.otf,
  UprightFont=*_Regular,
  ItalicFont=*_Italic,
  BoldFont=*_Bold,
  BoldItalicFont=*_Bold-Italic,
]

\setmonofont{Libertinus Mono}

%%% Para incluir o comando \url{}:

\usepackage[breaklinks=true]{hyperref}

%% Configurações ABNTeX:

\usepackage[hyphens,alf,abnt-and-type=e,abnt-last-names=bibtex,abnt-etal-cite=1,abnt-etal-list=1]{abntex2cite}


% Recuo no primeiro parágrafo:

\usepackage{indentfirst}

% Mudando recuo de parágrafo e tamanho da fonte:

\setlength{\parindent}{7mm}
\fontsize{11pt}{13.2pt}\selectfont

% Usado para evitar linhas órfãs:

\usepackage{needspace}

% Mudar espaçamento entre número e título de seção:

%% \def\l@figure{\@dottedtocline{1}{1.5em}{2.5em}}

\usepackage{tocloft}
\setlength{\cftfignumwidth}{2.55em} 

% Para escrever "Apêndices" e "Anexos":

%\usepackage[titletoc]{appendix}

% Elimina recuo nas notas de rodapé:

\usepackage[hang,flushmargin]{footmisc}

% Primeira página vazia em cada capítulo:

\makeatletter
\renewcommand\chapter{\if@openright\cleardoublepage\else\clearpage\fi
                    \thispagestyle{empty}% original style: plain
                    \global\@topnum\z@
                    \@afterindentfalse
                    \secdef\@chapter\@schapter}
\makeatother

% Diminui chance de linhas órfãs e viúvas:

\clubpenalty1000000
\widowpenalty1000000

% Comando para mudar a fonte de citação literal de código (verbatim):

\makeatletter
\newcommand{\verbatimfont}[1]{\def\verbatim@font{#1}}%
\makeatother

% Pacote para referir a capítulos pelo nome:
\usepackage{nameref}
\makeatletter
\newcommand*{\currentname}{\@currentlabelname}
\makeatother

% Definições (alternativas a nameref) para se referir a capítulos e seções pelo nome:
\let\Chaptermark\chaptermark
\def\chaptermark#1{\def\Chaptername{#1}\Chaptermark{#1}}
\let\Sectionmark\sectionmark
\def\sectionmark#1{\def\Sectionname{#1}\Sectionmark{#1}}
\let\Subsectionmark\subsectionmark
\def\subsectionmark#1{\def\Subsectionname{#1}\Subsectionmark{#1}}
\let\Subsubsectionmark\subsubsectionmark
\def\subsubsectionmark#1{\def\Subsubsectionname{#1}\Subsubsectionmark{#1}}

% Pacote para usar "clearpage" sem finalizar página, apenas descarregar todos os elementos já inseridos:

\usepackage{afterpage}

% Pacote para que páginas em branco não mostrem cabeçalhos nem rodapés:

\usepackage{emptypage}

% Pacote para índice Remissivo (a ser impresso por \printindex):

\usepackage{makeidx}
\makeindex

% Pacotes auxiliares no design da capa:

\usepackage{xcoffins}

%% Pacotes para usar o pacote Metapost para fazer graficos e diagramas:

\usepackage{luamplib}
\everymplib{input mpcolornames; beginfig(1);}
\everyendmplib{endfig;}

%% Pacotes para diagramas, desenhos, gráficos:

\usepackage{tikz}
\usetikzlibrary{backgrounds,patterns}

\usepackage{pgfplots}
\pgfplotsset{compat=1.12}

%% Pacote para ferramentas gráficas:

\usepackage{graphics}
\usepackage{graphicx} % Sem esse, alguns includegraphics não funcionam.

\usepackage{threeparttable} % Para alinhar figuras e legendas com "measuredfigure".

\usepackage{wrapfig} % Figuras ao lado do texto.

\usepackage{picinpar}% Alteranativa para figuras ao lado de texto. http://ctan.org/pkg/picinpar

\usepackage[font={small}, margin=0cm, justification=centering]{caption} % Legendas com fonte pequena.

% Comando para adicionar fonte de figuras e semelhantes:

\newcommand{\source}[1]{\captionsetup{singlelinecheck=false,justification=justified}\caption*{\footnotesize \noindent Fonte: {#1}}}
% Tamanhos de fontes matemáticas em relação à fonte do texto:
%% {tamanho do texto} {matemática} {matemática script} {matemática scriptscript}

\DeclareMathSizes{10}{10}{6}{4}
\DeclareMathSizes{9}{9}{5}{3}

\usepackage{amsthm} % deve ser chamado antes de mdframed conforme dito em http://tex.stackexchange.com/questions/283763/why-dont-i-get-non-italic-normal-font-inside-theorem-environment-using-newmdth

% Seleção de fontes de equações:

\usepackage[math-style=french]{unicode-math} % não funciona se compilar com pdflatex, deve-se usar xelatex ou lualatex

\usepackage{yfonts} % para usar caracteres góticos para algumas variaveis

% Pacote para maior controle de fluxo de figuras (usado para opção [H] que fixa a posição das figuras):

\usepackage{float}

% Pacote usado para impedir elementos flutuantes (figuras, tabelas, etc.) de aparecerem em seção errada:

\usepackage[section]{placeins}

% Pacote inserido para fazer o símbolo de grau:

\usepackage{gensymb}

% Pacote inserido para fazer símbolos de circuito elétrico:

\usepackage{marvosym}

% Definição de variáveis e unidades:

\newcommand{\angstrom}{\text{\normalfont\AA}}
\newcommand{\pol}{\ensuremath{pol}}
\newcommand{\cm}{\ensuremath{cm}}
\newcommand{\km}{\ensuremath{km}}
\newcommand{\hm}{\ensuremath{hm}}
\newcommand{\dam}{\ensuremath{dam}}
\newcommand{\dm}{\ensuremath{dm}}
\newcommand{\mol}{\ensuremath{mol}}
\newcommand{\mi}{\ensuremath{mi}}
\newcommand{\h}{\ensuremath{h}}
\newcommand{\s}{\ensuremath{s}}
\newcommand{\Min}{\ensuremath{min}}
\newcommand{\Pa}{\ensuremath{Pa}}
\newcommand{\atm}{\ensuremath{atm}}
\newcommand{\mmHg}{\ensuremath{mmHg}}
\newcommand{\BarP}{\ensuremath{Bar}}
\newcommand{\PsiP}{\ensuremath{PSI}}
\newcommand{\lb}{\ensuremath{lb}}
\newcommand{\N}{\ensuremath{N}}
\newcommand{\kg}{\ensuremath{kg}}
\newcommand{\kgf}{\ensuremath{kgf}}
\newcommand{\are}{\ensuremath{a}}
\newcommand{\litro}{\ensuremath{\ell}}
\newcommand{\g}{\ensuremath{g}}

% Símbolos de diferencial:

\def\D{\mathrm{d}} %% diferencial - comando \D{}
\def\DI{{\delta}} % diferencial inexata ``delta''

\newcommand*\diff{\mathop{}\!\mathrm{d}}
\newcommand*\Diff[1]{\mathop{}\!\mathrm{d^#1}}

% Evita que equações extrapolem a margem (e permite outros recursos tipográficos avançados):

\usepackage{microtype}

%% Simbolos matemáticos extra:

\DeclareMathSymbol{\Omega}{\mathalpha}{letters}{"0A}
\DeclareMathSymbol{\varOmega}{\mathalpha}{operators}{"0A}
\providecommand*{\upOmega}{\varOmega} % for siunitx

%% Pacote para uso de unidades SI com distância padronizada entre valor e unidade:

\usepackage{siunitx}
\sisetup{%
      binary-units=true,
      group-separator={.},
      group-digits=integer,
      load-configurations=abbreviations,
      load=addn,
      per-mode=fraction,
      output-decimal-marker={,},
      range-phrase= --,
      separate-uncertainty=true,
      math-ohm = \ensuremath{\upOmega}, % senão \ohm não funciona
      text-ohm = Ω,  % senão \ohm não funciona
    }

\DeclareSIUnit\milha{mi}
\DeclareSIUnit\polegada{pol}
\DeclareSIUnit\alqueire{alqueire}
\DeclareSIUnit\inch{in}
\DeclareSIUnit\foot{ft}
\DeclareSIUnit\kgf{kgf}
\DeclareSIUnit\lbf{lbf}
\DeclareSIUnit\pound{lb}
\DeclareSIUnit\rev{rev}
\DeclareSIUnit\rpm{rpm}
\DeclareSIUnit\pkt{PKT}
\DeclareSIUnit{\calorie}{cal}
\DeclareSIUnit{\cal}{cal}
\DeclareSIUnit{\Cal}{Cal}
\DeclareSIUnit{\Calorie}{\kilo\calorie}
\DeclareSIUnit{\fahrenheit}{\degree F}

\DeclareSIUnit{\nothing}{\relax} % usado para mostrar prefixos sem unidade

% Equações químicas:

\usepackage[version=4]{mhchem}

\usepackage{minted}
\usepackage{xcolor} % to access the named colour LightGray
\definecolor{LightGray}{gray}{0.9}

% Definição de ambientes exercício, solução, definição, exemplo, demonstração, etc.

\usepackage{mdframed}

\usepackage{environ} % usado nas definições \MakeDiscussionTopSecret, \MakeDiscussionsPublic, etc.

% Exercício e solução (pacote Exsheets):

\usepackage{exsheets} %[2015/02/09]

% Configurações do pacote Exsheets:

\SetupExSheets{
  headings          = block-subtitle ,
  subtitle-format   = \sc ,
  counter-within    = {chapter} , % Contar dentro dos capítulos.
  counter-format    = ch.qu[1] , % Formato 1.1, 1.2, 1.3, etc.
  label-format      = ch.qu[1] , % Formato 1.1, 1.2, 1.3, etc.
  headings-format   = \bfseries ,
  question/pre-hook = \needspace{0.3\textheight}\vspace{1ex} ,
  question/post-hook = \vspace{1ex} ,
  question/pre-body-hook = {\vspace{1ex} \mdframed[backgroundcolor=verde_UnB!10, linewidth=1pt, innermargin=+0cm, outermargin=+0cm]},
  question/post-body-hook = \endmdframed ,
  solution/pre-hook = \needspace{0.3\textheight}\vspace{1ex},
  solution/post-hook = \vspace{1ex},
  solution/pre-body-hook = \vspace{1ex},
  solution/print    = false
}

\renewcommand\thequestion{\thechapter.\arabic{question}} % Formato 1.1, 1.2, 1.3, etc., ao usar \ref{}:

% Ambiente "Definição":

\mdfdefinestyle{definitionSty}{backgroundcolor=gray!10, linecolor=verde_UnB, linewidth=0pt, innerleftmargin=3ex, innerrightmargin=3ex, innertopmargin=1ex, innerbottommargin=1ex, innermargin =0, outermargin =0, needspace={0.2\textheight}}
\newcounter{definitionCounter}[chapter]
\numberwithin{definitionCounter}{chapter}
\theoremstyle{definition}
\newmdtheoremenv[style=definitionSty]{definition}[definitionCounter]{Defini\c{c}\~{a}o}

% Ambiente "Teorema":

\mdfdefinestyle{theoremSty}{backgroundcolor=verde_UnB!10, linecolor=verde_UnB, linewidth=0pt, innerleftmargin=3ex, innerrightmargin=3ex, innertopmargin=1ex, innerbottommargin=1ex, innermargin=0, outermargin=0, needspace={0.2\textheight}}
\newcounter{theoremCounter}[chapter] 
\numberwithin{theoremCounter}{chapter}
\newmdtheoremenv[style=theoremSty]{theorem}[theoremCounter]{Teorema}

% Ambiente "Demonstração":

\mdfdefinestyle{demonstrationSty}{backgroundcolor=gray!10, linecolor=verde_UnB, linewidth=0pt, innerleftmargin=3ex, innerrightmargin=3ex, innertopmargin=1ex, innerbottommargin=1ex, innermargin =0cm, outermargin =0cm, needspace={0.2\textheight}}
\newcounter{demonstrationCounter}[chapter]
\numberwithin{demonstrationCounter}{chapter}
\newmdtheoremenv[style=demonstrationSty]{demonstration}[demonstrationCounter]{Demonstra\c{c}\~{a}o}

% Ambiente "Axioma":

\mdfdefinestyle{axiomSty}{backgroundcolor=verde_UnB!10, linecolor=verde_UnB, linewidth=1pt, innerleftmargin=3ex, innerrightmargin=3ex, innertopmargin=1ex, innerbottommargin=1ex, innermargin=0, outermargin=0, needspace={0.1\textheight}}
\newcounter{axiomCounter}[chapter] 
\numberwithin{axiomCounter}{chapter}
\newmdtheoremenv[style=axiomSty]{axiom}[axiomCounter]{Axioma}

% Ambiente "Exemplo":

\mdfdefinestyle{exampleSty}{backgroundcolor=gray!10, linecolor=verde_UnB, linewidth=0pt, innerleftmargin=3ex, innerrightmargin=3ex, innertopmargin=1ex, innerbottommargin=1ex,  innermargin =+0cm, outermargin =+0cm, needspace={0.2\textheight}}
\newcounter{exampleCounter}[chapter]
\numberwithin{exampleCounter}{chapter}
\newmdtheoremenv[style=exampleSty]{example}[exampleCounter]{Exemplo}

% Ambiente "Problema":

\mdfdefinestyle{problemSty}{backgroundcolor=gray!10, linecolor=black, linewidth=1pt, innerleftmargin=3ex, innerrightmargin=3ex, innertopmargin=1ex, innerbottommargin=1ex, innermargin =+0cm, outermargin =+0cm, needspace={0.15\textheight}}
\newcounter{problemCounter}[chapter]
\numberwithin{problemCounter}{chapter}
\newmdtheoremenv[style=problemSty]{problem}[problemCounter]{Problema}

% Ambiente "Nota":

\mdfdefinestyle{noteSty}{backgroundcolor=white, font=\small, fontcolor=verde_UnB, linecolor=verde_UnB, linewidth=1pt, innerleftmargin=3ex, innerrightmargin=3ex, innertopmargin=2ex, innerbottommargin=2ex, innermargin =+0cm, outermargin =+0cm, needspace={0.1\textheight}}
\newcounter{noteCounter}[chapter]
\numberwithin{noteCounter}{chapter}
\newmdtheoremenv[style=noteSty]{note}[noteCounter]{Nota}

% Ambiente "Discussão" (texto complementar):

\mdfdefinestyle{discussionSty}{backgroundcolor=white, fontcolor=verde_UnB, linecolor=verde_UnB, linewidth=0pt, innerleftmargin=3ex, innerrightmargin=3ex, innertopmargin=3ex, innerbottommargin=3ex, innermargin=+0cm, outermargin =+0cm, needspace={0.25\textheight}}
\newcounter{discussionCounter}%[chapter]
%\numberwithin{discussionCounter}{chapter}
\newmdtheoremenv[style=discussionSty]{discussion}[discussionCounter]{Texto Complementar}

% Ambiente "Dedicatória" (baseado em https://tex.stackexchange.com/a/167529/91816):

\newenvironment{dedication}
  {
    \itshape             % the text is in italics
    \raggedleft          % flush to the right margin
  }
  {
    \par % end the paragraph
  }
  % space at bottom is three times that at the top
  
% Citação com autor:

\renewenvironment{quotation}
  {\small\list{}{\rightmargin=0.1cm \leftmargin=4cm \topsep=1.5\baselineskip}%
   \item\relax}
  {\endlist}

\def\signed #1{{\leavevmode\unskip\nobreak\newline
  \hbox{}\nobreak #1 %
  \hfil\parfillskip=0pt \finalhyphendemerits=0 \endgraf}}

\newsavebox\mybox
\newenvironment{citacao}[1]
  {\savebox\mybox{#1}\begin{quotation}}
  {\signed{\usebox\mybox}\end{quotation}}


\usepackage{tcolorbox}

\title{Modelo de livro da série \\ Ensino de Graduação da Editora UnB}
\author{Leo Fernandes at al} % Para mais de um autor, use: \author{autor1 \\ autor2 \\ autor3}
\date{Agosto 2024}

% Cabeçalhos das páginas ímpares (odd), à esquerda (LO), centro (CO) e direita (RO):

\makeatletter
\fancyhead[LO]{}
\fancyhead[CO]{\small \if@mainmatter \small \leftmark \fi}
\fancyhead[RO]{}
\makeatother

% Cabeçalhos das páginas pares (even):
% Obs.: caso haja mais de um autor, substitua \theauthor pelo nome do 1º autor, e coloque os outros em CE ou RE:

\fancyhead[LE]{\theauthor}
\fancyhead[CE]{}
\fancyhead[RE]{IFAL} 

% Rodapé:

\fancyfoot[C]{\small \thepage}
%\fancyfoot[RO]{\small \nouppercase \rightmark}

% Para que \mainmatter não resete a contagem de páginas (https://tex.stackexchange.com/questions/380816/prevent-frontmatter-to-reset-counting):

\makeatletter
\renewcommand\mainmatter{%
  \cleardoublepage
  \@mainmattertrue
  \renewcommand{\thepage}{\arabic{page}}
  %\pagenumbering{arabic}% Don't reset
}
\makeatother

% Início do documento:

\begin{document}

% Ajustando espaços no sumário:

\makeatletter
\renewcommand{\@pnumwidth}{2.0em}
\renewcommand{\@tocrmarg}{2.7em}
\makeatother

% Parte Pré-Textual:

\frontmatter
%\pagestyle{empty}

% O código do texto está nos arquivos dentro da pasta "txt".
% Por exemplo, o comando %%% Física para Ciências Agrárias

%%% Capa

%%% Definindo estilo desta pagina como "vazio": 
\thispagestyle{empty}

%%% Capa: 
%%% Modelo tirado de http://tex.stackexchange.com/questions/17579/how-can-i-design-a-book-cover

\vspace*{\fill}

\begin{center}
\textbf{\color{verde_UnB}\fontsize{36pt}{45.6pt}\selectfont \textbf{Programação Web II}}
\textbf{\color{cinza_UnB}\fontsize{24pt}{28.6pt}\selectfont \textbf{Material de Apoio da Disciplina PWEB II do IFAL/Maceió}}

\end{center}

\vfill

\vspace*{\fill}

\clearpage inclui o código que está no arquivo "titulo.tex" dentro da pasta txt.
% Crie seus próprios arquivos à vontade e inclua-os aqui com o comando \include{}.

%%% Física para Ciências Agrárias

%%% Capa

%%% Definindo estilo desta pagina como "vazio": 
\thispagestyle{empty}

%%% Capa: 
%%% Modelo tirado de http://tex.stackexchange.com/questions/17579/how-can-i-design-a-book-cover

\vspace*{\fill}

\begin{center}
\textbf{\color{verde_UnB}\fontsize{36pt}{45.6pt}\selectfont \textbf{Programação Web II}}
\textbf{\color{cinza_UnB}\fontsize{24pt}{28.6pt}\selectfont \textbf{Material de Apoio da Disciplina PWEB II do IFAL/Maceió}}

\end{center}

\vfill

\vspace*{\fill}

\clearpage
\clearpage

% \thispagestyle{empty}

\begin{center}

\begin{tabular}{r:l}
						& \textbf{\Large Instituto Federal de Alagoas}				\\
						&														\\
	{\bf Reitor} 		& XXXXXXXXX									\\
	{\bf Vice-Reitor}	& XXXXXXXXXXXXX										\\
						&														\\
						& \textbf{\Large Editora YYYY}							\\
						&														\\
	{\bf Diretor}		& XXXXXXXXXX							\\
						&														\\
	{\bf Conselho editorial}	&	XXXXXXXXXXXXXXXXXXX					\\
								&	XXXXXXXXXXXXXXXXXXX				\\
								&	XXXXXXXXXXXXXXXXXXX			\\
								&	XXXXXXXXXXXXXXXXXXX	\\
								&	XXXXXXXXXXXXXXXXXXXX						\\
\end{tabular}

\end{center}

\clearpage

% \clearpage

% 1ª página de cada capítulo tem estilo "vazio" (sem cabeçalho ou rodapé):
\thispagestyle{empty}

\begin{center}

\vspace*{\fill}

\textbf{\color{verde_UnB}\fontsize{30pt}{35pt}\selectfont \textbf{Programação Web II}}

\vfill

{\textbf{\fontsize{13pt}{15.6pt}\selectfont \centering  Leo Fernandes, PhD}}

%\vspace{1cm}

%{\textbf{\fontsize{13pt}{15.6pt}\selectfont \centering  Nome do segundo autor}}

\vspace*{\fill}

\end{center}

\newpage
\clearpage

% \thispagestyle{empty}

\begin{center}

\begin{tabular}{r:l}
						& \textbf{\Large Universidade de Brasília}				\\
						&														\\
	{\bf Preparação e revisão} 		& Tiago de Aguiar Rodrigues					\\
	{\bf Diagramação}				& Wladimir de Andrade Oliveira				\\
						&														\\
						& \textcopyright 2018 Editora Universidade de Brasília	\\
						&														\\
						& Direitos exclusivos para esta edição:					\\
						&	Editora Universidade de Brasília					\\
						&	SCS, quadra 2, bloco C, nº 78, edifício OK,			\\
						&	2º andar, CEP 70302-4200							\\
						&	Site: \url{www.editora.unb.br}						\\
						&	E-mail: \href{mailto:contatoeditora@unb.br}{contatoeditora@unb.br}	\\							&																	\\
						& Todos os direitos reservados. \\
						& Nenhuma parte desta publicação	\\														& poderá ser armazenada ou reproduzida \\
						& por qualquer meio sem a autorização \\
						& por escrito da Editora. \\
						&		\\
						& {\bf Edital Livros Didáticos}	\\
						& Esta obra foi publicada com recursos provenientes	\\
						& do Edital DEG/UnB nº 13/2017.
\end{tabular}

\end{center}

\newpage

% \clearpage

% \thispagestyle{empty}

\vspace*{\fill}

\begin{dedication}
    A minha avozinha Renilda, \\
    que me ensinou a bordar, tocar violão, \\
    cuidar de peixes de aquário, produzir licor \\
    e montar um orquidário, \\
    e de quem herdei o diletantismo.    \vspace{\baselineskip}
%    \usefont{T1}{LobsterTwo-LF}{bx}{it}
%    Leonardo
\end{dedication}

% Para inserir várias dedicatórias, descomente as duas linhas acima para colocar a assinatura do autor e use o modelo abaixo para fazer uma outra dedicatória:

%\vspace{5\baselineskip}

%\begin{dedication}
%	
%    \vspace{\baselineskip}
%    \usefont{T1}{LobsterTwo-LF}{bx}{it}
%    Autor 2
%\end{dedication}
% \cleardoublepage

% \chapter*{Agradecimentos}
\thispagestyle{empty}

\hfill %\hspace{0.5cm}

...
% \cleardoublepage

% Lista de símbolos/variáveis:

% \chapter*{Lista de símbolos}
\thispagestyle{empty}

\vspace{0.5cm}

\begin{longtable}{|c|l|} % "|" = "linha vertical", "c" = "centro", "l" = "esquerda", "r" = "direita"
 \hline
  \textbf{Variável} & \textbf{Grandeza representada} \\ \hline
  $A$ & área \\ \hline
  $\alpha$, $\beta$, $\gamma$ & radiações \\ \hline
  $a$ & aceleração \\ \hline
  $\mathcal{B}$ & campo magnético \\ \hline
  $D$ & diâmetro \\ \hline
  $d$ & deslocamento, distância \\ \hline
  $\textgoth{D}$ & dose absorvida \\ \hline
  $\Delta{X}$ & variação ou margem de erro da medida $X$ \\ \hline
  $\D{X}$ & elemento infinitesimal, diferencial da grandeza $X$ \\ \hline
  $\DI{X}$ & diferencial inexata da grandeza $X$ \\ \hline
  $E$ & energia \\ \hline
  $\mathcal{E}$ & campo elétrico \\ \hline
  $\epsilon$  & permissividade elétrica \\ \hline
  $\eta$ & coeficiente de viscosidade \\ \hline
  $F$ & força \\ \hline
  $\phi$   & ângulo, ângulo azimutal (coordenadas polares ou esféricas) \\ \hline
  $h$ & altura, elevação \\ \hline
  $H$ & dose equivalente \\ \hline
  $I$ & corrente elétrica \\ \hline
  $K$ & energia cinética \\ \hline
  $k$ & constante elástica \\ \hline
  $k$ & número de onda, vetor de onda ($\vec{k}$) \\ \hline
  $\kappa$ & constante eletrostática de um meio \\ \hline
  $\kappa_0$ & constante de Coulomb (constante eletrostática do vácuo) \\ \hline
  $L$ & lado \\ \hline
  $L$ & momento angular \\ \hline
  $l$ & comprimento \\ \hline
  $m$ & massa \\ \hline
  $\mu$ & permeabilidade magnética \\ \hline
  $\mu_{din}$ & coeficiente de atrito dinâmico \\ \hline
  $\mu_{est}$ & coeficiente de atrito estático \\ \hline
  $N$ & força normal \\ \hline
  $\hat{n}$ & vetor normal unitário \\ \hline
  $P$ & peso \\ \hline
  $\mathscr{P}$ & potência \\ \hline
  $p$ & momento linear \\ \hline
  $\mathscr{p}$ & pressão \\ \hline
  $Q$ & calor \\ \hline
  $\mathscr{Q}$ & vazão \\ \hline
  $\textgoth{Q}$ & fator de qualidade \\ \hline  
  $q$ & carga elétrica \\ \hline
  $r$ & coordenada radial \\ \hline
  $r$, $\theta$ & coordenadas polares \\ \hline
  $r$, $\theta$, $\phi$ & coordenadas esféricas \\ \hline  
  $R$ & raio \\ \hline
  $\mathscr{R}$ & resistência elétrica \\ \hline
  $\rho$ & densidade \\ \hline
  $T$ & temperatura \\ \hline
  $\tau$ & torque, momento de força \\ \hline
  $\omega$ & velocidade angular \\ \hline
  $\mathscr{T}$ & força de tensão \\ \hline
  $T$ & temperatura \\ \hline
  $U$ & energia potencial \\ \hline
  $\mathscr{U}$ & diferença de potencial elétrico (voltagem) \\ \hline
  $v$ & velocidade \\ \hline
  $\mathscr{W}$ & trabalho \\ \hline 
  $x$ & deformação linear \\ \hline
  $\bar{x}$ & média de x \\ \hline
  $\vec{x}$ & vetor x \\ \hline
  $\mathscr{x}$, $\mathscr{y}$, $\mathscr{z}$ & eixos de um plano cartesiano \\ \hline
  $\textgoth{X}$ & exposição \\ \hline
\end{longtable}
% \cleardoublepage

% Lista de figuras:

% \listoffigures \thispagestyle{empty}
% \cleardoublepage

% Lista de tabelas:

% \listoftables \thispagestyle{empty}
% \cleardoublepage

% Sumário:

\tableofcontents\thispagestyle{empty}
\cleardoublepage 

% Parte Textual:

\mainmatter

% \chapter{Introdução aos Protocolos Web}
% \section{Protocolo de Transferência de Hipertexto (HTTP)}

% \subsubsection{Requisições e Respostas}
% \subsubsection{Estrutura das Requisições HTTP}
% \subsubsection{Estrutura das Respostas HTTP}
% \subsubsection{Cabeçalhos}
% \subsubsection{Métodos HTTP}
% \subsubsection{Códigos de Status}
% \subsubsection{Tipos MIME}
% \subsubsection{Cookies}
% \subsubsection{Executando HTTP na Linha de Comando}
% \begin{itemize}
% \item Usando cURL e outras ferramentas
% \end{itemize}


% \section{Comunicação Bidirecional}
% \subsubsection{Polling e Long Polling}
% \subsubsection{Server-Sent Events}
% \subsubsection{WebSockets}
% \begin{itemize}
% \item Introdução ao WebSockets
% \item Implementando WebSockets no Node.js
% \end{itemize}


\chapter{A Web}

\section{História da Web}

A história da Web é uma jornada fascinante que começou no final da década de 1980 e se desenvolveu rapidamente, transformando a maneira como interagimos com a informação e o mundo ao nosso redor.

% \subsection{As Origens da Web}
A World Wide Web foi inventada por Tim Berners-Lee, um cientista britânico do CERN (Organização Europeia para a Pesquisa Nuclear), em 1989. Berners-Lee desenvolveu a Web como uma solução para compartilhar documentos entre os pesquisadores do CERN. Ele propôs um sistema de gerenciamento de informações que conectaria documentos utilizando hiperlinks, permitindo que os usuários navegassem facilmente de um documento para outro .

Em 1990, Berners-Lee, juntamente com o engenheiro de sistemas Robert Cailliau, desenvolveu os primeiros componentes essenciais da Web:
\begin{itemize}
    \item \textbf{HTML (HyperText Markup Language):} A linguagem de marcação usada para criar documentos na Web.
    \item \textbf{HTTP (Hypertext Transfer Protocol):} O protocolo usado para a transferência de documentos pela Web.
    \item \textbf{URL (Uniform Resource Locator):} O sistema de endereçamento que permite identificar e acessar documentos na Web.
\end{itemize}

No final de 1990, Tim Berners-Lee tinha o primeiro servidor Web e navegador instalado e funcionando no CERN, demonstrando suas ideias. Ele desenvolveu o código para seu servidor Web em um computador NeXT. Para evitar que fosse desligado acidentalmente, o computador trazia uma etiqueta escrita à mão em tinta vermelha: "This machine is a server. DO NOT POWER IT DOWN!!"

O endereço do servidor web era \textit{info.cern.ch} \url{http://info.cern.ch/}, foi lançado em 1990 no CERN, e a primeira página web (ainda ativa em \url{http://info.cern.ch/hypertext/WWW/TheProject.html}) explicou o que era a World Wide Web e como utilizá-la.

% \subsection{A Popularização da Web}
A Web começou a se popularizar em 1993, com o lançamento do navegador Mosaic, desenvolvido por Marc Andreessen e Eric Bina no National Center for Supercomputing Applications (NCSA). O Mosaic foi o primeiro navegador a suportar gráficos embutidos, além de texto, em uma interface amigável. Isso marcou o início da explosão da Web como uma ferramenta acessível para o público em geral.

Em 1994, Berners-Lee fundou o World Wide Web Consortium (W3C)~\footnote{ O World Wide Web Consortium (W3C) desenvolve padrões e diretrizes para ajudar todos a construir uma web baseada nos princípios de acessibilidade, internacionalização, privacidade e segurança. - \url{https://www.w3.org/}} no MIT, com o objetivo de desenvolver padrões abertos que garantissem o crescimento contínuo da Web de maneira estável e interoperável. O W3C continua a ser a principal organização de padrões para a Web até hoje .

% \subsection{Evolução e Impacto da Web}
Desde o seu início, a Web passou por várias evoluções importantes. A transição de uma Web estática, onde as páginas eram basicamente documentos de leitura, para uma Web dinâmica e interativa, onde aplicações complexas podem ser executadas diretamente no navegador, foi um marco significativo. Tecnologias como JavaScript, CSS e AJAX contribuíram para essa evolução.

A Web também teve um impacto profundo na sociedade, revolucionando a comunicação, o comércio, a educação, o entretenimento e praticamente todos os aspectos da vida moderna. Com o surgimento da Web 2.0, a interação social online tornou-se predominante, dando origem a plataformas como redes sociais, blogs e wikis.

Hoje, a Web continua a evoluir, com tendências emergentes como a Web Semântica, a Web3 e a integração de tecnologias de Inteligência Artificial, prometendo transformar ainda mais a maneira como vivemos e trabalhamos.

Para mais informações sobre a história da web, acesse ~\cite{cern_history_web}.


% \begin{tcolorbox}[colback=gray!10, colframe=gray!80, title=Internet x Web]
% É comum que os termos "Internet" e "Web" sejam usados como sinônimos, mas eles referem-se a conceitos diferentes. A Internet é uma infraestrutura global de redes interconectadas que permite a comunicação entre computadores ao redor do mundo. Ela é composta por cabos, roteadores, servidores e protocolos que facilitam o tráfego de dados entre dispositivos. Por outro lado, a Web, ou World Wide Web, é a especificação de um conjunto de tecnologias que opera sobre a Internet, utilizando o protocolo HTTP para acessar e interagir com documentos e recursos disponíveis em servidores ao redor do mundo. A Web é apenas uma das muitas aplicações que utilizam a Internet, ao lado de outras como e-mail, FTP e VoIP. Em resumo, a Internet é a rede que conecta os computadores, enquanto a Web é uma maneira de acessar e compartilhar informações através dessa rede.
% \end{tcolorbox}

\begin{question}
    Qual é a diferença entre Internet e Web?
\end{question}

\begin{solution}
    É comum que os termos "Internet" e "Web" sejam usados como sinônimos, mas eles referem-se a conceitos diferentes. A Internet é uma infraestrutura global de redes interconectadas que permite a comunicação entre computadores ao redor do mundo. Ela é composta por cabos, roteadores, servidores e protocolos que facilitam o tráfego de dados entre dispositivos. Por outro lado, a Web, ou World Wide Web, é a especificação de um conjunto de tecnologias que opera sobre a Internet, utilizando o protocolo HTTP para acessar e interagir com documentos e recursos disponíveis em servidores ao redor do mundo. A Web é apenas uma das muitas aplicações que utilizam a Internet, ao lado de outras como e-mail, FTP e VoIP. Em resumo, a Internet é a rede que conecta os computadores, enquanto a Web é uma maneira de acessar e compartilhar informações através dessa rede.
\end{solution}




\section{Protocolo de Transferência de Hipertexto (HTTP)}

\subsection{Requisições e Respostas}
O Protocolo de Transferência de Hipertexto (HTTP) é a base da comunicação na web. Em essência, o HTTP funciona no modelo de requisição-resposta, onde um cliente (geralmente um navegador) faz uma requisição a um servidor, e o servidor responde com os dados solicitados.

O cliente envia uma requisição HTTP que consiste em uma linha de requisição, cabeçalhos e um corpo opcional. O servidor responde com uma resposta HTTP, que também contém uma linha de status, cabeçalhos e um corpo opcional.

% Sugestão de Figura: Uma figura mostrando o fluxo de uma requisição HTTP e a respectiva resposta.

\subsection{Estrutura das Requisições HTTP}
Uma requisição HTTP possui a seguinte estrutura básica:

\begin{itemize}
    \item \textbf{Linha de Requisição:} Contém o método HTTP, o caminho da URL, e a versão do protocolo HTTP.
    \item \textbf{Cabeçalhos:} Metadados que fornecem informações adicionais sobre a requisição. Exemplo: `Host`, `User-Agent`, `Accept`.
    \item \textbf{Corpo (opcional):} Contém os dados enviados com a requisição, como dados de formulários em uma requisição POST.
\end{itemize}

\begin{verbatim}
GET /index.html HTTP/1.1
Host: www.example.com
User-Agent: Mozilla/5.0
Accept: text/html

\end{verbatim}

\subsection{Estrutura das Respostas HTTP}
A resposta HTTP segue uma estrutura similar:

\begin{itemize}
    \item \textbf{Linha de Status:} Inclui a versão do protocolo, um código de status, e uma frase descritiva.
    \item \textbf{Cabeçalhos:} Informações sobre o servidor e os dados que estão sendo enviados. Exemplo: `Content-Type`, `Content-Length`, `Set-Cookie`.
    \item \textbf{Corpo (opcional):} Contém o conteúdo real da resposta, como uma página HTML.
\end{itemize}

\begin{verbatim}
HTTP/1.1 200 OK
Content-Type: text/html
Content-Length: 1234

<html>
  <body>...</body>
</html>
\end{verbatim}

\subsection{Cabeçalhos}
Os cabeçalhos HTTP são essenciais para o controle de requisições e respostas. Eles podem ser divididos em quatro categorias principais:

\begin{itemize}
    \item \textbf{Cabeçalhos Gerais:} Aplicam-se tanto a requisições quanto a respostas, como `Cache-Control` e `Date`.
    \item \textbf{Cabeçalhos de Requisição:} Específicos para requisições, como `Accept`, `User-Agent`, e `Host`.
    \item \textbf{Cabeçalhos de Resposta:} Específicos para respostas, como `Content-Type`, `Server`, e `Set-Cookie`.
    \item \textbf{Cabeçalhos de Entidade:} Relacionam-se ao corpo da mensagem, como `Content-Length` e `Content-Encoding`.
\end{itemize}

% Sugestão de Tabela: Tabela listando alguns cabeçalhos comuns e suas funções.

\subsection{Métodos HTTP}
Os métodos HTTP definem a ação a ser realizada na requisição. Os métodos mais comuns são:

\begin{itemize}
    \item \textbf{GET:} Solicita a representação de um recurso específico. Não deve alterar o estado do servidor.
    \item \textbf{POST:} Envia dados para o servidor para criar ou modificar um recurso.
    \item \textbf{PUT:} Substitui o recurso no servidor pela carga útil da requisição.
    \item \textbf{DELETE:} Remove um recurso específico do servidor.
    \item \textbf{HEAD:} Solicita uma resposta idêntica a GET, mas sem o corpo de resposta.
    \item \textbf{PATCH:} Aplica modificações parciais a um recurso.
\end{itemize}

\subsection{Códigos de Status}
Os códigos de status HTTP informam o resultado da requisição. Eles são divididos em cinco classes:

\begin{itemize}
    \item \textbf{1xx (Informativo):} Indica que a requisição foi recebida e está em processamento.
    \item \textbf{2xx (Sucesso):} Indica que a requisição foi bem-sucedida. Exemplo: 200 OK.
    \item \textbf{3xx (Redirecionamento):} Indica que o cliente precisa tomar uma ação adicional para completar a requisição. Exemplo: 301 Moved Permanently.
    \item \textbf{4xx (Erro do Cliente):} Indica que houve um erro na requisição. Exemplo: 404 Not Found.
    \item \textbf{5xx (Erro do Servidor):} Indica que o servidor encontrou um erro ao processar a requisição. Exemplo: 500 Internal Server Error.
\end{itemize}

\subsection{Tipos MIME}
O tipo MIME (Multipurpose Internet Mail Extensions) define o tipo de conteúdo da resposta HTTP. Os tipos MIME comuns incluem:

\begin{itemize}
    \item \textbf{text/html:} Para documentos HTML.
    \item \textbf{application/json:} Para dados no formato JSON.
    \item \textbf{image/png:} Para imagens no formato PNG.
    \item \textbf{application/xml:} Para dados no formato XML.
\end{itemize}

% Sugestão de Tabela: Tabela com exemplos de tipos MIME comuns e suas descrições.

\subsection{Cookies}
Cookies são pequenos pedaços de dados armazenados no cliente pelo servidor. Eles são usados para manter o estado entre requisições HTTP, que são inerentemente sem estado. Cookies podem armazenar informações como identificadores de sessão, preferências do usuário, e outros dados de rastreamento.

Os cookies são definidos no cabeçalho `Set-Cookie` da resposta HTTP e são enviados de volta ao servidor em requisições subsequentes.

% Sugestão de Figura: Diagrama ilustrando o fluxo de cookies entre o cliente e o servidor.

\subsection{Executando HTTP na Linha de Comando}
Ferramentas como cURL permitem que desenvolvedores façam requisições HTTP diretamente da linha de comando. Isso é útil para testar APIs e depurar problemas de rede.

\begin{verbatim}
# Exemplo de requisição GET com cURL
curl -X GET https://api.example.com/data

# Exemplo de requisição POST com cURL
curl -X POST https://api.example.com/data -d '{"key":"value"}' -H "Content-Type: application/json"
\end{verbatim}

% Comentário: Incluir mais exemplos práticos de cURL ou outras ferramentas como HTTPie pode ser útil para os alunos.


\section{A linguagem de marcação de hyper-texto (HTML)}

A \textbf{Linguagem de Marcação de Hyper-Texto (HTML)} é a espinha dorsal da Web, sendo a principal linguagem utilizada para criar e estruturar conteúdo na Internet. HTML é uma linguagem de marcação, o que significa que ela usa uma série de \textit{tags} ou marcas para definir elementos dentro de um documento. Esses elementos podem incluir textos, imagens, links, tabelas, formulários, e muito mais.

\subsection{Estrutura Básica de um Documento HTML}

Todo documento HTML possui uma estrutura básica que deve ser seguida. A estrutura de um arquivo HTML típico é composta por:

\begin{verbatim}
<!DOCTYPE html>
<html>
  <head>
    <title>Título da Página</title>
  </head>
  <body>
    <h1>Bem-vindo à Web</h1>
    <p>Este é um exemplo de documento HTML.</p>
  </body>
</html>
\end{verbatim}

\begin{itemize}
    \item \textbf{DOCTYPE:} Declara o tipo do documento e a versão do HTML que está sendo utilizada. O mais comum atualmente é o \texttt{<!DOCTYPE html>} que indica HTML5.
    \item \textbf{<html>:} A tag \texttt{<html>} encapsula todo o conteúdo do documento HTML.
    \item \textbf{<head>:} A tag \texttt{<head>} contém metadados sobre o documento, como o título que aparece na aba do navegador, links para folhas de estilo, scripts, e outras informações.
    \item \textbf{<body>:} A tag \texttt{<body>} contém todo o conteúdo visível para o usuário, como textos, imagens, links, e outros elementos de interface.
\end{itemize}

\subsection{Elementos HTML}
Os elementos HTML são representados por tags, que geralmente vêm em pares: uma tag de abertura e uma tag de fechamento. As tags de abertura e fechamento delimitam o início e o fim de um elemento.

\begin{itemize}
    \item \textbf{Tags de Cabeçalho:} Usadas para criar títulos e subtítulos na página (\texttt{<h1>} a \texttt{<h6>}).
    \item \textbf{Tags de Parágrafo:} Usadas para definir parágrafos de texto (\texttt{<p>}).
    \item \textbf{Tags de Link:} Usadas para criar hiperlinks para outros documentos ou seções (\texttt{<a href="...">}).
    \item \textbf{Tags de Imagem:} Usadas para incluir imagens no documento (\texttt{<img src="...">}).
    \item \textbf{Tags de Lista:} Usadas para criar listas ordenadas (\texttt{<ol>}) ou não ordenadas (\texttt{<ul>}).
\end{itemize}

\subsection{A Evolução do HTML}

O HTML evoluiu significativamente desde sua criação. As primeiras versões do HTML focavam em estruturação básica de conteúdo, mas versões posteriores, como o HTML5, introduziram novos elementos semânticos, como \texttt{<article>}, \texttt{<section>}, e \texttt{<nav>}, além de suporte nativo para multimídia e APIs avançadas para desenvolvimento web.

O HTML5, em particular, tornou-se o padrão mais amplamente adotado, sendo projetado para ser compatível com dispositivos modernos e suportando uma gama mais ampla de funcionalidades, como vídeos embutidos, gráficos, e aplicativos interativos.

\section{O sistema de endereçamento (URL)}

A \textbf{Uniform Resource Locator (URL)}, ou Localizador Uniforme de Recursos, é um sistema de endereçamento utilizado na Web para identificar e localizar recursos disponíveis na Internet. Cada URL aponta para um recurso específico, que pode ser uma página HTML, um documento PDF, uma imagem, ou qualquer outro tipo de conteúdo acessível na Web.

\subsection{Estrutura de uma URL}
Uma URL é composta por várias partes distintas que indicam diferentes aspectos do recurso a ser acessado. A estrutura geral de uma URL é a seguinte:

\scriptsize
\begin{verbatim}
http://www.exemplo.com:80/caminho/para/recurso?query=parametro#ancora
\end{verbatim}
\normalsize

Vamos decompor essa URL em suas partes componentes:

\begin{itemize}
    \item \textbf{Protocolo:} \texttt{http://} ou \texttt{https://} especifica o protocolo de comunicação a ser utilizado. O protocolo HTTPS é uma versão segura do HTTP.
    \item \textbf{Nome de Domínio:} \texttt{www.exemplo.com} identifica o servidor onde o recurso está hospedado. Este nome de domínio pode ser traduzido em um endereço IP pelo DNS (Domain Name System).
    \item \textbf{Porta:} \texttt{:80} é a porta de comunicação do servidor utilizada pelo protocolo. A porta 80 é o padrão para HTTP, enquanto a porta 443 é padrão para HTTPS.
    \item \textbf{Caminho:} \texttt{/caminho/para/recurso} indica o caminho no servidor para o recurso específico. Pode referenciar diretórios ou arquivos.
    \item \textbf{Query String:} \texttt{?query=parametro} é uma cadeia de caracteres que fornece parâmetros adicionais para a requisição, frequentemente utilizados para busca ou filtragem de dados.
    \item \textbf{Fragmento:} \texttt{\#ancora} refere-se a uma parte específica do recurso, como uma âncora em uma página HTML que leva a um local específico dentro do documento.
\end{itemize}

\subsection{Funcionalidade das URLs}

URLs desempenham um papel crucial na navegação na Web. Elas não apenas indicam ao navegador onde encontrar um recurso, mas também como acessar e interagir com ele. Por exemplo, uma URL pode incluir parâmetros que modificam a forma como o conteúdo é apresentado ou como uma pesquisa é realizada em um site.

\subsection{URLs e SEO}

As URLs também são um elemento importante na otimização para motores de busca (SEO). URLs claras, descritivas e bem estruturadas podem melhorar a indexação de uma página pelos motores de busca e aumentar a visibilidade do site.

\subsection{Segurança e URLs}

Com o crescimento da Web, a segurança das URLs tornou-se uma preocupação central. HTTPS, por exemplo, foi introduzido para proteger as comunicações entre o navegador e o servidor, criptografando os dados para evitar interceptações maliciosas.

% Comentário: Considerar a inclusão de exemplos práticos que mostrem como diferentes tipos de URLs são construídos e utilizados em aplicações web.




\section{Comunicação Bidirecional}
\subsection{Polling e Long Polling}
Polling e Long Polling são técnicas usadas para manter uma comunicação bidirecional entre o cliente e o servidor, especialmente em aplicações que necessitam de atualizações frequentes.

\begin{itemize}
    \item \textbf{Polling:} O cliente faz requisições periódicas ao servidor para verificar se há novas informações. Essa técnica pode ser ineficiente devido à quantidade de requisições feitas.
    \item \textbf{Long Polling:} O servidor mantém a conexão aberta até que haja novas informações para enviar ao cliente, reduzindo o número de requisições.
\end{itemize}

\subsection{Server-Sent Events}
Server-Sent Events (SSE) permitem que um servidor envie dados de atualização para o cliente de forma contínua, usando uma conexão HTTP única. Isso é útil para notificações em tempo real, feeds de dados e outras aplicações que necessitam de atualizações contínuas.

\begin{verbatim}
# Exemplo de código SSE no lado do servidor (Node.js)
const express = require('express');
const app = express();

app.get('/events', (req, res) => {
    res.setHeader('Content-Type', 'text/event-stream');
    res.setHeader('Cache-Control', 'no-cache');
    res.setHeader('Connection', 'keep-alive');

    setInterval(() => {
        res.write(`data: ${new Date().toISOString()}\n\n`);
    }, 1000);
});

app.listen(3000, () => console.log('SSE server running on port 3000'));
\end{verbatim}

\subsection{WebSockets}
\textbf{Introdução ao WebSockets}

WebSockets fornecem uma comunicação bidirecional em tempo real entre o cliente e o servidor sobre uma única conexão TCP. Diferente do HTTP, o WebSocket permite que tanto o cliente quanto o servidor enviem dados a qualquer momento, sem a necessidade de repetidas requisições HTTP.

\textbf{Implementando WebSockets no Node.js}

Para implementar WebSockets em Node.js, a biblioteca \texttt{ws} é uma escolha popular. A seguir, um exemplo básico de implementação:

\begin{verbatim}
const WebSocket = require('ws');

const wss = new WebSocket.Server({ port: 8080 });

wss.on('connection', ws => {
    ws.on('message', message => {
        console.log(`Received message => ${message}`);
    });
    ws.send('Hello! You are connected.');
});
\end{verbatim}

Este exemplo cria um servidor WebSocket que ouve na porta 8080. Quando um cliente se conecta, o servidor pode enviar e receber mensagens em tempo real.

% Comentário: Considere incluir uma explicação sobre como usar ferramentas como o "WebSocket Tester" ou extensões de navegador para testar WebSockets.















% =================================================================================
% \chapter{Usando APIs Web}
% \section{Alterando Páginas Web Dinamicamente Usando a API DOM}
% \begin{itemize}
%     \item O Modelo de Objeto de Documento (DOM)
%     \item Os Diferentes Tipos de Nós
%     \item Selecionando Elementos
%     \item Modificando Elementos
%     \item Criando, Adicionando e Excluindo Elementos
%     \item Exemplo Prático: Criação Dinâmica de uma Tabela
% \end{itemize}
% \section{Carregando Dados Sincronamente via Ajax e a API Fetch}
% \begin{itemize}
%     \item Comunicação Sincrona versus Assíncrona
%     \item Carregando Dados via Ajax
% \end{itemize}
% \subsection{Carregando Dados via a API Fetch}
% \begin{itemize}
%     \item Visão Geral das APIs Web
%     \item Suporte dos Navegadores para APIs Web
% \end{itemize}



\chapter{Usando APIs Web}

\section{Alterando Páginas Web Dinamicamente Usando a API DOM}

\subsection{O Modelo de Objeto de Documento (DOM)}
O Modelo de Objeto de Documento, ou DOM (\textit{Document Object Model}), é uma interface de programação que permite a manipulação de documentos HTML e XML como uma estrutura de árvore, onde cada nó representa uma parte do documento, como elementos, atributos ou texto. O DOM é essencial para a criação de páginas web dinâmicas, permitindo que os desenvolvedores interajam e modifiquem a estrutura e o conteúdo das páginas em tempo real.

% Sugestão de Figura: Uma ilustração mostrando a estrutura em árvore do DOM, com elementos HTML como nós.

\subsection{Os Diferentes Tipos de Nós}
No DOM, existem vários tipos de nós, cada um representando uma parte específica do documento:

\begin{itemize}
    \item \textbf{Nós de Elemento:} Representam tags HTML, como \texttt{<div>}, \texttt{<p>}, \texttt{<a>}.
    \item \textbf{Nós de Atributo:} Representam os atributos de um elemento HTML, como \texttt{id}, \texttt{class}.
    \item \textbf{Nós de Texto:} Representam o conteúdo textual dentro de um elemento HTML.
    \item \textbf{Nós de Comentário:} Representam os comentários no código HTML.
    \item \textbf{Nós de Documento:} Representam o próprio documento inteiro, como a raiz da árvore DOM.
\end{itemize}

% Sugestão de Código: Exemplo de como acessar diferentes tipos de nós no DOM usando JavaScript.

\begin{verbatim}
// Acessando um nó de elemento
let elementNode = document.getElementById('myElement');

// Acessando um nó de atributo
let attributeNode = elementNode.getAttributeNode('class');

// Acessando um nó de texto
let textNode = elementNode.firstChild;
\end{verbatim}

\subsection{Selecionando Elementos}
Para modificar uma página web dinamicamente, primeiro precisamos selecionar os elementos que desejamos alterar. O DOM fornece várias maneiras de selecionar elementos:

\begin{itemize}
    \item \textbf{getElementById}: Seleciona um único elemento pelo seu \texttt{id}.
    \item \textbf{getElementsByClassName}: Seleciona todos os elementos que possuem uma determinada classe.
    \item \textbf{getElementsByTagName}: Seleciona todos os elementos com uma determinada tag.
    \item \textbf{querySelector}: Seleciona o primeiro elemento que corresponde a um seletor CSS.
    \item \textbf{querySelectorAll}: Seleciona todos os elementos que correspondem a um seletor CSS.
\end{itemize}

\begin{verbatim}
// Exemplo de seleção de elementos
let element = document.getElementById('header');
let elements = document.querySelectorAll('.item');
\end{verbatim}

\subsection{Modificando Elementos}
Depois de selecionar um elemento, podemos modificar suas propriedades, conteúdo, ou estilo diretamente através do DOM.

\begin{itemize}
    \item \textbf{Modificando o Conteúdo:} Use \texttt{innerHTML} ou \texttt{textContent} para alterar o conteúdo interno de um elemento.
    \item \textbf{Alterando Estilos:} Modifique o estilo de um elemento através da propriedade \texttt{style}.
    \item \textbf{Adicionando/Removendo Classes:} Use \texttt{classList} para adicionar, remover ou alternar classes de um elemento.
\end{itemize}

\begin{verbatim}
// Exemplo de modificação de conteúdo e estilo
let element = document.getElementById('header');
element.textContent = 'Novo Título';
element.style.color = 'blue';
element.classList.add('novo-estilo');
\end{verbatim}

\subsection{Criando, Adicionando e Excluindo Elementos}
O DOM permite que você crie novos elementos, adicione-os ao documento e remova elementos existentes.

\begin{itemize}
    \item \textbf{Criando Elementos:} Use \texttt{document.createElement()} para criar um novo elemento.
    \item \textbf{Adicionando Elementos:} Use \texttt{appendChild()} ou \texttt{insertBefore()} para adicionar o novo elemento à árvore DOM.
    \item \textbf{Excluindo Elementos:} Use \texttt{removeChild()} ou \texttt{element.remove()} para remover um elemento.
\end{itemize}

\begin{verbatim}
// Exemplo de criação e adição de um novo elemento
let newElement = document.createElement('p');
newElement.textContent = 'Este é um novo parágrafo.';
document.body.appendChild(newElement);

// Exemplo de remoção de um elemento
let elementToRemove = document.getElementById('oldElement');
elementToRemove.remove();
\end{verbatim}

\subsection{Exemplo Prático: Criação Dinâmica de uma Tabela}
Vamos criar dinamicamente uma tabela HTML utilizando JavaScript e a API DOM. Este exemplo prático ajudará a solidificar os conceitos discutidos.

\begin{verbatim}
// Criando a tabela e os elementos associados
let table = document.createElement('table');
let headerRow = document.createElement('tr');
let headers = ['Nome', 'Idade', 'Profissão'];

headers.forEach(headerText => {
    let header = document.createElement('th');
    header.textContent = headerText;
    headerRow.appendChild(header);
});

table.appendChild(headerRow);

// Adicionando linhas de dados à tabela
let data = [
    ['Alice', 30, 'Engenheira de Software'],
    ['Bob', 25, 'Designer'],
    ['Carol', 28, 'Gerente de Projetos']
];

data.forEach(rowData => {
    let row = document.createElement('tr');
    rowData.forEach(cellData => {
        let cell = document.createElement('td');
        cell.textContent = cellData;
        row.appendChild(cell);
    });
    table.appendChild(row);
});

// Adicionando a tabela ao corpo do documento
document.body.appendChild(table);
\end{verbatim}

% Sugestão de Figura: Captura de tela da tabela gerada dinamicamente no navegador.

\section{Carregando Dados Sincronamente via Ajax e a API Fetch}

\subsection{Comunicação Sincrona versus Assíncrona}
A comunicação síncrona envolve a execução de tarefas de maneira sequencial, onde uma tarefa deve ser concluída antes que a próxima comece. Em contrapartida, a comunicação assíncrona permite que as tarefas sejam iniciadas sem esperar pela conclusão das anteriores, o que é crucial para melhorar a performance e a experiência do usuário em aplicações web.

% Comentário: Explicar o impacto da comunicação síncrona em termos de bloqueio da interface do usuário e como a comunicação assíncrona mitiga esse problema.

\subsection{Carregando Dados via Ajax}
Ajax (Asynchronous JavaScript and XML) é uma técnica amplamente usada para carregar dados de forma assíncrona sem recarregar a página inteira. Ajax permite que as aplicações web se tornem mais interativas e responsivas.

\begin{verbatim}
// Exemplo básico de uma requisição Ajax usando XMLHttpRequest
let xhr = new XMLHttpRequest();
xhr.open('GET', 'https://api.example.com/data', true);

xhr.onload = function() {
    if (xhr.status === 200) {
        console.log(xhr.responseText);
    } else {
        console.error('Erro na requisição.');
    }
};

xhr.send();
\end{verbatim}

% Sugestão de Comentário: Discutir a obsolescência do Ajax em favor do Fetch API, destacando a sintaxe mais moderna e as vantagens do Fetch.

\subsection{Carregando Dados via a API Fetch}

\subsubsection{Visão Geral das APIs Web}
A API Fetch é uma API moderna e baseada em Promises para fazer requisições HTTP assíncronas. Diferente do XMLHttpRequest, Fetch proporciona uma maneira mais simples e poderosa de buscar recursos através da rede.

\begin{verbatim}
// Exemplo básico de uma requisição com a Fetch API
fetch('https://api.example.com/data')
    .then(response => {
        if (!response.ok) {
            throw new Error('Erro na requisição');
        }
        return response.json();
    })
    .then(data => console.log(data))
    .catch(error => console.error('Erro:', error));
\end{verbatim}

\subsubsection{Suporte dos Navegadores para APIs Web}
Embora a API Fetch seja amplamente suportada pelos navegadores modernos, é importante considerar o suporte para navegadores legados ao desenvolver aplicações web. O uso de polyfills pode ser uma solução para garantir compatibilidade.

% Comentário: Incluir referências para documentação e ferramentas de verificação de suporte a navegadores, como o "Can I Use".


















% =================================================================================
\chapter{JavaScript no Lado do Servidor com Node.js}

\section{Introdução ao Node.js}
\subsubsection{Arquitetura do Node.js}
\begin{itemize}
\item Arquitetura Orientada a Eventos no Node.js
\item I/O Não Bloqueante
\end{itemize}
\subsubsection{Um Primeiro Programa}
\begin{itemize}
\item Construindo um Servidor Web Simples
\end{itemize}
\subsubsection{Gerenciamento de Pacotes com npm}
\begin{itemize}
\item Instalando e Gerenciando Pacotes
\item Criando e Publicando seu Próprio Pacote
\end{itemize}

\section{Utilizando Módulos Nativos do Node.js}
\subsubsection{Trabalhando com o Sistema de Arquivos}
\begin{itemize}
\item Lendo Arquivos
\item Escrevendo Arquivos
\item Deletando Arquivos
\end{itemize}
\subsubsection{Trabalhando com Buffers e Streams}
\subsubsection{Rede}
\begin{itemize}
\item Criando Servidores TCP e HTTP
\item Trabalhando com Sockets
\end{itemize}


\section{Construindo Servidores Web}
\subsubsection{Preparando o Ambiente de Desenvolvimento}
\subsubsection{Fornecendo Arquivos Estáticos}
\subsubsection{Usando o Framework Express.js}
\begin{itemize}
\item Configurando o Express.js
\item Roteamento no Express.js
\item Middleware no Express.js
\end{itemize}
\subsubsection{Processando Dados de Formulários}
\subsubsection{Usando Motores de Template com Express.js}
\begin{itemize}
\item Pug, Mustache, e EJS
\end{itemize}
\subsection{Extra: Formatos Web}
\begin{itemize}
    \item Formatos de dados: CSV, XML e JSON
    \item Formatos de imagem: JPG, GIF, PNG, SVG, WebP
    \item Comprando formatos de imagem
    \item Programas para processar imagens
    \item Formatos de audio e vídeo 
\end{itemize}













% =================================================================================

\chapter{Implementando Serviços Web}
\section{Introdução aos Serviços Web}

\subsubsection{Visão Geral dos Serviços Web}


\section{Serviços Web RESTful}

\subsubsection{Princípios do REST}
\subsubsection{Implementando uma API REST}
\begin{itemize}
\item Usando Express.js para Criar Endpoints RESTful
\item Manipulando Operações CRUD
\end{itemize}
\subsubsection{Consumindo uma API REST}
\begin{itemize}
\item Consumindo APIs RESTful com JavaScript
\item Tratamento de Erros e Melhores Práticas
\end{itemize}


\section{GraphQL}

\subsubsection{As Desvantagens do REST}
\subsubsection{Introdução ao GraphQL}
\subsubsection{Implementando GraphQL com Node.js}
\begin{itemize}
\item Configurando o GraphQL no Node.js
\item Escrevendo Queries e Mutations
\item GraphQL vs REST
\end{itemize}












% =================================================================================


\chapter{Armazenando Dados em Bancos de Dados}
\section{Bancos de Dados Relacionais}

\subsubsection{Entendendo Bancos de Dados Relacionais}
\subsubsection{A Linguagem SQL}
\begin{itemize}
\item Consultas SQL Básicas (SELECT, INSERT, UPDATE, DELETE)
\item Joins, Subconsultas e Transações
\end{itemize}
\subsubsection{Usando MySQL com Node.js}
\begin{itemize}
\item Configurando MySQL no Node.js
\item Realizando Operações CRUD com MySQL
\end{itemize}
\subsubsection{Mapeamento Objeto-Relacional (ORM)}
\begin{itemize}
\item Usando Sequelize com Node.js
\end{itemize}


\section{Bancos de Dados Não-Relacionais}

\subsubsection{Bancos de Dados Relacionais vs Não-Relacionais}
\subsubsection{Visão Geral dos Bancos de Dados Não-Relacionais}
\begin{itemize}
\item Bancos de Dados de Chave-Valor
\item Bancos de Dados Orientados a Documentos
\item Bancos de Dados de Grafos
\item Bancos de Dados Colunares
\end{itemize}
\subsubsection{Usando MongoDB com Node.js}
\begin{itemize}
\item Configurando MongoDB no Node.js
\item Operações CRUD com MongoDB
\item Framework de Agregação no MongoDB
\end{itemize}
\subsubsection{Comparando MySQL e MongoDB}




























% =================================================================================

\chapter{Arquiteturas Web}
\section{Arquiteturas em Camadas}

\subsubsection{Estrutura Básica das Arquiteturas em Camadas}
\subsubsection{Arquitetura Cliente-Servidor (Arquitetura de Dois Níveis)}
\subsubsection{Arquitetura de Múltiplos Níveis}
\begin{itemize}
\item Arquitetura N-Camadas
\item Vantagens e Desvantagens
\item Other Web APIs: 
\end{itemize}


\section{Arquiteturas Monolíticas e Distribuídas}

\subsubsection{Arquitetura Monolítica}
\subsubsection{Arquitetura Orientada a Serviços (SOA)}
\subsubsection{Arquitetura de Microsserviços}
\begin{itemize}
\item Comparando SOA e Microsserviços
\item Implementando Microsserviços com Node.js
\end{itemize}
\subsubsection{Arquitetura Baseada em Componentes}
\subsubsection{Arquitetura de Microfrontends}
\subsubsection{Arquitetura de Mensagens}
\begin{itemize}
\item Arquitetura Orientada a Eventos
\item Filas de Mensagens e Brokers (RabbitMQ, Kafka)
\end{itemize}
\subsubsection{Arquitetura de Serviços Web}
\begin{itemize}
\item SOAP vs REST
\item Implementando Serviços Web com Node.js
\end{itemize}


\section{Arquiteturas MV*}

\subsubsection{Model-View-Controller (MVC)}
\begin{itemize}
\item Implementando MVC com Express.js
\end{itemize}
\subsubsection{Model-View-Presenter (MVP)}
\subsubsection{Model-View-ViewModel (MVVM)}












% =================================================================================

\chapter{Padrões de Projeto}
\section{Introdução aos Padrões de Design}
\subsubsection{O que são Padrões de Design?}
\begin{itemize}
\item Importância dos Padrões de Design em Desenvolvimento de Software
\item Categorias de Padrões de Design: Criacionais, Estruturais e Comportamentais
\end{itemize}

\section{Padrões Criacionais}

\subsection{Singleton}
\subsection{Factory Method}
\subsection{Abstract Factory}
\subsection{Builder}
\subsection{Prototype}
\begin{itemize}
\item Implementando Padrões Criacionais com JavaScript/Node.js
\end{itemize}


\section{Padrões Estruturais}

\subsection{Adapter}
\subsection{Decorator}
\subsection{Facade}
\subsection{Proxy}
\subsection{Composite}
\subsection{Flyweight}
\subsection{Bridge}
\begin{itemize}
\item Implementando Padrões Estruturais com JavaScript/Node.js
\end{itemize}


\section{Padrões Comportamentais}

\subsection{Observer}
\subsection{Strategy}
\subsection{Command}
\subsection{Chain of Responsibility}
\subsection{Template Method}
\subsection{State}
\subsection{Iterator}
\subsection{Mediator}
\begin{itemize}
\item Implementando Padrões Comportamentais com JavaScript/Node.js
\end{itemize}















% =================================================================================

\chapter{Testando Aplicações Web}
\section{Introdução aos Testes}
\subsubsection{A Importância dos Testes no Desenvolvimento Web}

\section{Testes Unitários}
\subsubsection{Escrevendo Testes Unitários em JavaScript}
\begin{itemize}
\item Usando Mocha e Chai
\item Desenvolvimento Orientado por Testes (TDD)
\end{itemize}


\section{Testes de Integração}

\subsubsection{Testando APIs RESTful}
\begin{itemize}
\item Usando Supertest com Mocha
\end{itemize}


\section{Testes de Ponta a Ponta (E2E)}

\subsubsection{Automatizando Testes de Ponta a Ponta}
\begin{itemize}
\item Usando Cypress para Testes E2E
\end{itemize}













% =================================================================================

\chapter{Implantação e Hospedagem de Aplicações Web}
\section{Preparando para Implantação}

\subsubsection{Processo de Build e Release}
\subsubsection{Variáveis de Ambiente e Configurações}


\section{Implantando Aplicações}

\subsubsection{Implantando em Provedores de Nuvem}
\begin{itemize}
\item AWS, Azure e Google Cloud Platform
\end{itemize}
\subsubsection{Usando Docker para Implantação}
\begin{itemize}
\item Containerização com Docker
\item Dockerizando Aplicações Node.js
\end{itemize}
\subsubsection{Usando Pipelines CI/CD}
\begin{itemize}
\item Automatizando a Implantação com Jenkins, GitHub Actions e GitLab CI
\end{itemize}















% =================================================================================

\chapter{Segurança em Aplicações Web}
\section{Compreendendo a Segurança Web}

\subsubsection{Ameaças Comuns à Segurança}
\begin{itemize}
\item Injeção de SQL
\item Cross-Site Scripting (XSS)
\item Cross-Site Request Forgery (CSRF)
\end{itemize}


\section{Implementando Segurança no Node.js}

\subsubsection{Usando Helmet para Cabeçalhos de Segurança}
\subsubsection{Protegendo Aplicações Express}
\begin{itemize}
\item Validação e Sanitização de Entrada
\item Limitação de Taxa e Proteção contra DDoS
\end{itemize}
\subsubsection{Autenticação e Autorização}
\begin{itemize}
\item Implementando OAuth e JWT
\end{itemize}















% =================================================================================

\chapter{Otimizando a Performance de Aplicações Web}
\section{Identificando Gargalos de Performance}

\subsubsection{Ferramentas de Monitoramento e Profiling}
\begin{itemize}
\item Usando PM2 e New Relic
\end{itemize}


\section{Técnicas de Otimização de Performance}

\subsubsection{Estratégias de Cache}
\begin{itemize}
\item Usando Redis para Cache
\end{itemize}
\subsubsection{Balanceamento de Carga}
\subsubsection{Otimizando Consultas ao Banco de Dados}
\begin{itemize}
\item Indexação e Otimização de Consultas
\end{itemize}
\subsubsection{Otimização de Código}
\begin{itemize}
\item Vazamentos de Memória e Coleta de Lixo no Node.js
\item Otimização de Programação Assíncrona e do Event Loop
\end{itemize}


% \chapter{Sobre o Node.js}

% \thispagestyle{empty} 

% \begin{mdframed}[style=noteSty]

% {\center \textsc{Texto motivador} \par}

%    Este é um breve texto motivador. Espero que você esteja motivado(a)!
   
% \end{mdframed}


% \begin{itemize}

% \begin{enumerate}
%     \item ...
%     \item ;...
%     \item ...
% \end{enumerate}
 
% \end{itemize}

% % Muitos cientistas gostam de usar \LaTeX porque essa ferramenta possibilita escrever facilmente equações como a seguinte:
% % \begin{equation}
% %  \mathscr{p}+\frac{1}{2}{\rho}v^2+{\rho}gh = \text{constante}
% %  \label{eq:Bernoulli}
% % \end{equation}
% % onde $\mathscr{p}$ é a pressão, $v$ é a velocidade e $h$ é a elevação, ou seja, a “altura do tubo”. Essa equação pode ser deduzida a partir do \textit{Teorema Trabalho-Energia}. \index{Teorema Trabalho-Energia}  % entrada para o índice remissivo

% \newpage

% \chapter{Introdução ao Node.js}

% O Node.js é uma plataforma de código aberto construída sobre o motor JavaScript V8 do Google Chrome, que permite executar JavaScript no lado do servidor. Foi criado por Ryan Dahl em 2009 e desde então tem ganhado uma enorme popularidade na comunidade de desenvolvimento de software.

% \section{História}

% A história do Node.js remonta a 2009, quando Ryan Dahl apresentou pela primeira vez esta tecnologia em uma conferência. Naquela época, o desenvolvimento de aplicativos web era predominantemente baseado em tecnologias como PHP, Ruby on Rails e Java. No entanto, essas tecnologias enfrentavam desafios de escalabilidade quando se tratava de lidar com conexões simultâneas de clientes.

% Dahl percebeu que o modelo de E/S (entrada/saída) bloqueante das tecnologias existentes era um gargalo para o desempenho e a escalabilidade dos aplicativos web. Ele propôs o uso do motor JavaScript V8 do Chrome, que é conhecido por sua rápida execução de código, para criar um ambiente de servidor assíncrono e orientado a eventos.

% Assim, o Node.js nasceu com a promessa de oferecer um ambiente de execução leve e eficiente para construir aplicativos de rede escaláveis e em tempo real.

% \section{Por que é popular?}

% O Node.js ganhou popularidade por várias razões:

% \begin{enumerate}
%     \subsubsection{Arquitetura orientada a eventos e não bloqueante:} O Node.js utiliza uma arquitetura orientada a eventos que permite lidar com várias conexões de forma eficiente e escalável, sem bloquear a execução do código. Isso é particularmente útil para aplicativos que precisam lidar com um grande número de operações de E/S, como servidores web e aplicativos em tempo real.
    
%     \subsubsection{Ecossistema robusto:} O Node.js possui um vasto ecossistema de módulos e bibliotecas disponíveis através do npm, o maior repositório de pacotes de software do mundo. Isso facilita a criação de aplicativos complexos ao permitir que os desenvolvedores reutilizem código existente e incorporem funcionalidades de terceiros em seus aplicativos.
    
%     \subsubsection{JavaScript em ambos os lados:} Com o Node.js, os desenvolvedores podem usar JavaScript tanto no lado do cliente quanto no lado do servidor, o que simplifica o desenvolvimento de aplicativos web completos usando uma linguagem de programação coesa em toda a pilha tecnológica.
    
%     \subsubsection{Desempenho:} O Node.js é conhecido por seu desempenho rápido devido à execução assíncrona de código e ao motor JavaScript V8 altamente otimizado.
% \end{enumerate}

% Em resumo, o Node.js se tornou popular devido à sua capacidade de criar aplicativos de rede escaláveis e em tempo real de forma eficiente, sua vasta comunidade de desenvolvedores e seu ecossistema robusto de módulos e bibliotecas. Ele continua sendo uma escolha popular para o desenvolvimento de aplicativos web modernos e altamente performáticos. 

% Agora, vamos exemplificar com um pequeno código que demonstra a execução assíncrona no Node.js:

% \begin{minted}[
% frame=single, 
% linenos,
% firstnumber=1,
% framesep=2mm,
% baselinestretch=1.2,
% bgcolor=LightGray,
% fontsize=\footnotesize
% ]{js}
% // Exemplo de código assíncrono no Node.js
% const fs = require('fs');

% // Lê um arquivo de forma assíncrona
% fs.readFile('arquivo.txt', 'utf8', (err, data) => {
%   if (err) {
%     console.error('Erro ao ler o arquivo:', err);
%     return;
%   }
%   console.log('Conteúdo do arquivo:', data);
% });

% console.log('Operação assíncrona iniciada. 
%        Este console.log é executado antes da leitura do arquivo.');
% \end{minted}

% Neste exemplo, o código \texttt{fs.readFile()} inicia uma operação de leitura de arquivo de forma assíncrona. Enquanto o Node.js está lendo o arquivo, ele continua executando o código restante, como o \texttt{console.log()} subsequente, sem esperar a conclusão da operação de leitura do arquivo. Isso demonstra como o Node.js é capaz de lidar com operações de entrada/saída de forma não bloqueante, permitindo que outras operações continuem enquanto espera uma resposta de E/S.


% \chapter{Principais Características do Node.js}

% O Node.js é conhecido por várias características que o tornam uma escolha popular para o desenvolvimento de aplicativos de rede em tempo real. Aqui estão algumas das principais características destacadas:

% \begin{enumerate}
%     \subsubsection{Arquitetura Orientada a Eventos:}
%         \begin{itemize}
%             \item O Node.js opera com uma arquitetura orientada a eventos, onde as operações de entrada e saída são tratadas de forma assíncrona e não bloqueante.
%             \item Isso significa que em vez de esperar por uma operação de entrada/saída (como leitura de arquivos ou consultas de banco de dados) ser concluída antes de continuar, o Node.js continua a executar outras tarefas e manipula o resultado da operação quando estiver disponível.
%             \item Essa abordagem permite que o Node.js lide com um grande número de conexões simultâneas sem a necessidade de threads adicionais, o que resulta em um consumo mais eficiente de recursos.
%         \end{itemize}
        
%     \subsubsection{Não Bloqueante e Assíncrono:}
%         \begin{itemize}
%             \item O Node.js é não bloqueante e assíncrono por natureza, o que significa que as operações de E/S (entrada/saída) não bloqueiam a execução do código.
%             \item Em vez disso, o Node.js usa chamadas de retorno (callbacks), promessas ou async/await para lidar com operações assíncronas, permitindo que outras partes do código continuem executando enquanto aguardam a conclusão das operações de E/S.
%             \item Essa abordagem é fundamental para criar aplicativos escaláveis e eficientes, especialmente em cenários de rede em tempo real, onde várias solicitações podem chegar simultaneamente.
%         \end{itemize}
        
%     \subsubsection{Adequado para Aplicativos de Rede em Tempo Real:}
%         \begin{itemize}
%             \item Devido à sua natureza orientada a eventos e não bloqueante, o Node.js é particularmente adequado para desenvolver aplicativos de rede em tempo real, como aplicativos de bate-papo, jogos multiplayer online, streaming de vídeo e áudio, entre outros.
%             \item A capacidade do Node.js de lidar com um grande número de conexões simultâneas de forma eficiente e escalável o torna uma escolha popular para esses tipos de aplicativos.
%         \end{itemize}
% \end{enumerate}

% \section{Exemplo Prático}

% \begin{minted}[
% frame=single, 
% linenos,
% firstnumber=1,
% framesep=2mm,
% baselinestretch=1.2,
% bgcolor=LightGray,
% fontsize=\footnotesize
% ]{js}
% const http = require('http');

% // Cria um servidor HTTP simples
% const server = http.createServer((req, res) => {
%   // Simula uma operação de E/S assíncrona
%   setTimeout(() => {
%     res.writeHead(200, {'Content-Type': 'text/plain'});
%     res.end('Olá, mundo!\n');
%   }, 1000); // Responde após 1 segundo
% });

% // Inicia o servidor na porta 3000
% server.listen(3000, 'localhost', () => {
%   console.log('Servidor Node.js rodando em http://localhost:3000/');
% });
% \end{minted}

% Neste exemplo, criamos um servidor HTTP simples usando o módulo \texttt{http} do Node.js. O servidor responde com "Olá, mundo!" após um atraso de 1 segundo, simulando uma operação de E/S assíncrona. Isso demonstra como o Node.js pode lidar com operações assíncronas de forma eficiente, permitindo que o servidor continue a atender outras solicitações enquanto aguarda a conclusão da operação de E/S.


% \chapter{Explorando o Ecossistema Node.js}

% \section{npm (Node Package Manager)}

% O npm é o gerenciador de pacotes padrão para o ecossistema Node.js. Ele permite que os desenvolvedores instalem, compartilhem e gerenciem dependências de seus projetos de forma fácil e eficiente. Alguns pontos importantes a serem destacados sobre o npm incluem:

% \begin{itemize}
%     \subsubsection{Instalação de Pacotes:} Como instalar pacotes npm em um projeto Node.js usando o comando \mintinline{bash}{npm install}.
%     \subsubsection{Gerenciamento de Dependências:} Como adicionar, atualizar e remover dependências do projeto usando o npm.
%     \subsubsection{Arquivo \texttt{package.json}:} O papel crucial do arquivo \texttt{package.json} no gerenciamento de dependências, scripts de construção e metadados do projeto.
%     \subsubsection{Scripts npm:} A capacidade de definir e executar scripts personalizados usando o npm, como scripts de construção, execução de testes e muito mais.
% \end{itemize}

% \section{Express.js}

% Express.js é um framework web minimalista e flexível para Node.js, projetado para facilitar a construção de aplicativos web e APIs. Abordaremos os seguintes aspectos do Express.js:

% \begin{itemize}
%     \subsubsection{Configuração Básica:} Como configurar um aplicativo Express.js e definir rotas básicas.
%     \subsubsection{Middleware:} O conceito de middleware no Express.js e como usar middleware para manipular solicitações HTTP.
%     \subsubsection{Rotas:} Como definir e manipular rotas em um aplicativo Express.js para lidar com diferentes tipos de solicitações HTTP.
%     \subsubsection{Renderização de Visualizações:} Integração de mecanismos de visualização, como EJS ou Handlebars, para renderizar visualizações dinâmicas em aplicativos Express.js.
%     \subsubsection{Gestão de Erros:} Estratégias para lidar com erros e exceções em aplicativos Express.js, incluindo middleware de tratamento de erros.
% \end{itemize}

% \section{Exemplo Prático}

% Instalação do Express.js via npm

% \begin{minted}[
% frame=single, 
% linenos,
% firstnumber=1,
% framesep=2mm,
% baselinestretch=1.2,
% bgcolor=LightGray,
% fontsize=\footnotesize
% ]{shell}
% $ npm install express
% \end{minted}

% Após a instalação via NPM, você pode criar o exemplo abaixo em um projeto node e depois executar como uma aplicação Node normal.

% \begin{minted}[
% frame=single, 
% linenos,
% firstnumber=1,
% framesep=2mm,
% baselinestretch=1.2,
% bgcolor=LightGray,
% fontsize=\footnotesize
% ]{js}


% // Arquivo index.js
% const express = require('express');
% const app = express();

% // Rota principal
% app.get('/', (req, res) => {
%   res.send('Bem-vindo ao meu aplicativo Express!');
% });

% // Rota de exemplo com parâmetros
% app.get('/user/:id', (req, res) => {
%   res.send(`ID do usuário: ${req.params.id}`);
% });

% // Middleware para tratamento de erros
% app.use((err, req, res, next) => {
%   console.error(err.stack);
%   res.status(500).send('Algo deu errado!');
% });

% // Iniciar o servidor na porta 3000
% app.listen(3000, () => {
%   console.log('Servidor Express em execução na porta 3000');
% });
% \end{minted}

% Este exemplo demonstra como criar um aplicativo básico usando o Express.js, definir rotas simples e usar um middleware para lidar com erros. O npm é utilizado para instalar o Express.js e outras dependências necessárias para o projeto.

\cleardoublepage

% \chapter{Figuras e gráficos}

\thispagestyle{empty} 

Sugiro que você guarde todas as figuras na pasta ``figs'' para que seu projeto fique mais organizado. A figura \ref{fig:logolatex} mostra como é fácil inserir uma figura com legenda e referência à fonte.

\begin{figure}
	\centering
	\begin{minipage}{0.6\linewidth}
		\centering
		\caption{Logo \LaTeX.}
		\label{fig:logolatex}
		\includegraphics[width=\linewidth]{figs/1280px-LaTeX-logo.png}
		\source{Wikimedia Commons \cite{wikimedia-latex}.}
	\end{minipage}
\end{figure}

Além de figuras, \index{figuras} é possível inserir caixas de texto de diversos tipos, como axiomas, teoremas etc. Elas podem ser configuradas e novos tipos delas podem ser criados no arquivo \begin{verbatim}config/envs.tex\end{verbatim}.

Existem pacotes que permitem criar figuras e gráficos no próprio código \LaTeX. Por exemplo, temos

\begin{itemize}
    \item PGFPlots \url{http://pgfplots.sourceforge.net/}
    \item TikZ \url{http://www.texample.net/tikz/examples/all/}
    \item Metapost \url{http://tex.loria.fr/prod-graph/zoonekynd/metapost/metapost.html}
    \item PSTricks \url{https://tug.org/PSTricks/main.cgi?file=examples}
\end{itemize}

\begin{question}
    Explique como Isaac Newton usaria cada um dos pacotes seguintes, se vivesse no tempo presente:
    \begin{enumerate}[label=(\Alph*)]
        \item Metapost
        \item TikZ
        \item PGFPlots
        \item PSTricks
    \end{enumerate}
\end{question}

\begin{solution}
    \begin{enumerate}[label=(\Alph*)]
        \item Para fazer figuras 3D.
        \item Para fazer diagramas.
        \item Para traçar gráficos.
        \item Para fazer de um tudo.
    \end{enumerate}
\end{solution}

\section{Soluções deste capítulo}

\printsolutions
% \clearpage
% %\fancyfoot[CO,CE]{}
\fancyfoot[LO]{\small}
\fancyfoot[RO]{\small}
%\fancyfoot[CE]{\Author}
\fancyfoot[LE]{\small}
\fancyfoot[RE]{\small}

\begin{discussion}

\hfill

\begin{center}
  \textsc{\large Criando quadros}\\
\end{center}

\vspace{1ex}

\begin{flushright}
  \textbf{Leonardo Luiz e Castro} \\
  \vspace{1ex}
  \textit{\small Instituto de Física} \\
  \textit{\small Universidade de Brasília} \\
\end{flushright}

\addcontentsline{toc}{section}{Texto complementar -   Criando quadros (Leonardo Luiz e Castro)}

\vspace{1ex}

\section*{O que são quadros?}

Quadros são parecidos com tabelas, com a diferença de que não contêm dados numéricos, mas sim informação textual. Aparentemente, a ABNT é pioneira no mundo em se preocupar com tal distinção.

\section*{Exemplo de quadro em \LaTeX?}

Preferi listar os quadros como figuras. Vi algum livro que fazia assim... Veja como a figura \ref{fig:estados-da-materia} ficou interessante! Lembre-se que só é possível inserir uma figura dentro de um ambiente como este se a opção de não flutuação ([H]) for utilizada.

\begin{figure}[H]
	\centering
	\begin{minipage}{\hsize}
		\centering
		\taburulecolor{white}\arrayrulecolor{white}
		\caption{Quadro de caracterização dos estados da matéria.}
		\label{fig:estados-da-materia}
		\begin{tabular}{| >{\centering\cellcolor{verde_UnB}\color{white}}p{0.11\hsize} | >{\centering\cellcolor{verde_UnB}\color{white}}p{0.12\hsize} | >{\centering\cellcolor{verde_UnB}\color{white}}p{0.12\hsize} | >{\centering\cellcolor{verde_UnB}\color{white}\footnotesize}p{0.11\hsize} | >{\centering\cellcolor{verde_UnB}\color{white}\footnotesize}p{0.16\hsize} | >{\centering\cellcolor{verde_UnB}\color{white}\footnotesize}p{0.17\hsize} |}
			\hline
			{\bf Estado} & {\bf Volume} & {\bf Forma} & {\bf É fluido?} & {\bf É matéria condensada?} & {\bf É compressível?} \tabularnewline
			\hline
		\end{tabular}
		\begin{tabular}{| >{\centering\cellcolor{verde_UnB!70}\color{white}\small}p{0.11\hsize} | >{\centering\cellcolor{verde_UnB!50}\color{black}\small}p{0.12\hsize} | >{\centering\cellcolor{verde_UnB!50}\color{black}\small}p{0.12\hsize} | >{\centering\cellcolor{verde_UnB!50}\color{black}\small}p{0.11\hsize} | >{\centering\cellcolor{verde_UnB!50}\color{black}\small}p{0.16\hsize} | >{\centering\cellcolor{verde_UnB!50}\color{black}\small}p{0.17\hsize} |}
			\hline
			Gasoso & ajustável & ajustável & sim & não & muito \tabularnewline
			\hline
			Líquido & fixo & ajustável & sim & sim & pouco \tabularnewline
			\hline
			Sólido & fixo & fixa & não & sim & não \tabularnewline
			\hline
		\end{tabular}
		\source{elaboração do autores de Física para Ciências Agrárias e Ambientais (Leonardo Castro e Olavo Filho).}
	\end{minipage}
\end{figure}

\end{discussion}

% \cleardoublepage

% \chapter{Ambientes}

\thispagestyle{empty}

Este modelo disponibiliza alguns ``ambientes'', ou seja, caixas de texto com formatação especial para certos tipos de elementos que são automaticamente numerados (e.g. teorema 1.1, teorema 1.2 etc.). Ambientes podem ser criados e configurados por edição do arquivo envs.tex da pasta config.

\section{Exemplos de ambientes disponíveis}

\begin{axiom}
    \LaTeX produz equações mais bonitas que qualquer editor WYSIWYG.
\end{axiom}

\begin{theorem}\textsc{teorema LaTeX-WYSIWYG}
    Todo físico prefere usar código \LaTeX puro que qualquer editor WYSIWYG.
\end{theorem}

\index{\LaTeX} % entrada para o índice remissivo
\index{WYSIWYG}

\begin{demonstration}
    Físicos gostam de equações bonitas. Editores What-You-See-Is-What-You-Get não são apropriados para fazer equações bonitas.\footnote{É certo que há editores WYSIWYG baseados em \LaTeX, mas eles não nos dão o mesmo nível de controle.} Logo, se algum físico preferisse usar um editor WYSIWYG no lugar de \LaTeX, não seria muito inteligente. Como todo físico é inteligente, o teorema está demonstrado \textit{ad absurdum}.
\end{demonstration}

\begin{example}
    Einstein usaria um editor WYSIWYG ou \LaTeX? \\
    Einstein era físico. Portanto, usando o teorema LaTeX-WYSIWYG, concluímos que ele usaria \LaTeX.
\end{example}

\begin{question}
    Einstein usaria um editor WYSIWYG ou \LaTeX?
\end{question}

\begin{solution}
    Einstein era físico. Portanto, usando o teorema LaTeX-WYSIWYG, concluímos que ele usaria \LaTeX.
\end{solution}

\begin{question}
    Marie Curie usaria um editor WYSIWYG ou \LaTeX?
\end{question}

\begin{solution}
    Deixamos esta sem resposta para o estudante se esforçar mais.
\end{solution}

\clearpage

\section{Soluções deste capítulo}

\printsolutions
% \cleardoublepage

% \chapter{Conclusão}

\thispagestyle{empty}

Você deve começar a editar o seu livro agora mesmo!
% \cleardoublepage

% Apêndices:

\renewcommand\appendixname{APÊNDICE}
\renewcommand\appendixpagename{APÊNDICES}
\renewcommand{\appendixtocname}{APÊNDICES}

\begin{appendices}
\chapter{O Ambiente de Desenvolvimento}

Este apêndice descreve as etapas necessárias para configurar o ambiente de desenvolvimento utilizado ao longo deste livro. O ambiente será configurado no sistema operacional Linux Ubuntu, utilizando o editor de código Visual Studio Code (VS Code), a linguagem Node.js, o framework Express, e os bancos de dados MySQL e MongoDB.

\section{Instalando o Ubuntu}
\begin{itemize}
    \item Certifique-se de que você tenha a versão mais recente do Ubuntu instalada em sua máquina. Você pode baixar a versão mais recente do Ubuntu a partir do site oficial: \url{https://ubuntu.com/download}.
    \item Siga as instruções para instalar o Ubuntu, configurando o particionamento do disco e outras opções conforme necessário.
    \item Após a instalação, execute o comando abaixo para garantir que todos os pacotes estejam atualizados:
    \begin{verbatim}
    sudo apt update && sudo apt upgrade -y
    \end{verbatim}
\end{itemize}

\section{Instalando o Visual Studio Code (VS Code)}
\begin{itemize}
    \item Baixe e instale o Visual Studio Code (VS Code), um editor de código leve, mas poderoso, ideal para o desenvolvimento em Node.js.
    \item Execute o seguinte comando para instalar o VS Code via linha de comando:
    \begin{verbatim}
    sudo snap install --classic code
    \end{verbatim}
    \item Após a instalação, abra o VS Code e configure as extensões recomendadas para desenvolvimento em Node.js:
    \begin{itemize}
        \item \textbf{Node.js Extension Pack}
        \item \textbf{Prettier - Code Formatter}
        \item \textbf{ESLint}
    \end{itemize}
\end{itemize}

\section{Instalando o Node.js e npm}
\begin{itemize}
    \item O Node.js é uma plataforma JavaScript que permite executar código JavaScript no lado do servidor.
    \item Para instalar o Node.js e o npm (Node Package Manager) em Ubuntu, use o seguinte comando:
    \begin{verbatim}
    sudo apt install nodejs npm
    \end{verbatim}
    \item Verifique a instalação executando:
    \begin{verbatim}
    node -v
    npm -v
    \end{verbatim}
    Esses comandos devem retornar as versões instaladas do Node.js e do npm.
\end{itemize}

\section{Instalando o Framework Express}
\begin{itemize}
    \item O Express é um framework web para Node.js, amplamente utilizado para construir aplicações web robustas e escaláveis.
    \item Para instalar o Express globalmente, execute:
    \begin{verbatim}
    npm install -g express-generator
    \end{verbatim}
    \item Para criar uma nova aplicação Express, use:
    \begin{verbatim}
    express nome-do-projeto
    \end{verbatim}
    \item Navegue até o diretório do projeto e instale as dependências:
    \begin{verbatim}
    cd nome-do-projeto
    npm install
    \end{verbatim}
\end{itemize}

\section{Instalando o MySQL}
\begin{itemize}
    \item O MySQL é um sistema de gerenciamento de banco de dados relacional que será utilizado para armazenar dados estruturados.
    \item Para instalar o MySQL, execute:
    \begin{verbatim}
    sudo apt install mysql-server
    \end{verbatim}
    \item Após a instalação, inicie o serviço MySQL e configure a segurança inicial:
    \begin{verbatim}
    sudo systemctl start mysql
    sudo mysql_secure_installation
    \end{verbatim}
    \item Conecte-se ao MySQL utilizando:
    \begin{verbatim}
    sudo mysql -u root -p
    \end{verbatim}
\end{itemize}

\section{Instalando o MongoDB}
\begin{itemize}
    \item O MongoDB é um banco de dados NoSQL orientado a documentos que será utilizado para armazenar dados não estruturados.
    \item Para instalar o MongoDB, execute os seguintes comandos:
    \begin{verbatim}
    sudo apt install -y mongodb
    \end{verbatim}
    \item Inicie o serviço MongoDB:
    \begin{verbatim}
    sudo systemctl start mongodb
    \end{verbatim}
    \item Verifique se o MongoDB está funcionando corretamente:
    \begin{verbatim}
    sudo systemctl status mongodb
    \end{verbatim}
    \item Para interagir com o MongoDB, utilize o cliente de linha de comando:
    \begin{verbatim}
    mongo
    \end{verbatim}
\end{itemize}

\section{Configurando o Ambiente de Desenvolvimento}
\begin{itemize}
    \item \textbf{Configuração do Ambiente Node.js no VS Code:}
    \begin{itemize}
        \item Abra o VS Code e carregue o projeto Node.js.
        \item Configure o arquivo \texttt{launch.json} no VS Code para depuração de aplicações Node.js.
    \end{itemize}
    
    \item \textbf{Conectando ao MySQL e MongoDB:}
    \begin{itemize}
        \item Para se conectar ao MySQL em uma aplicação Node.js, utilize o pacote \texttt{mysql}:
        \begin{verbatim}
        npm install mysql
        \end{verbatim}
        
        \item Para se conectar ao MongoDB, utilize o pacote \texttt{mongoose}:
        \begin{verbatim}
        npm install mongoose
        \end{verbatim}
    \end{itemize}
\end{itemize}

\section{Testando a Configuração}
\begin{itemize}
    \item Para garantir que tudo está configurado corretamente, crie um pequeno projeto Node.js que se conecta ao MySQL e ao MongoDB, e execute-o usando o VS Code.
    
    \item Verifique se o projeto pode iniciar o servidor Express e realizar operações básicas de CRUD (Create, Read, Update, Delete) nos bancos de dados MySQL e MongoDB.
    
    \item Abaixo está um exemplo básico de uma aplicação Node.js que se conecta ao MySQL:
    \begin{verbatim}
    const mysql = require('mysql');
    const connection = mysql.createConnection({
      host: 'localhost',
      user: 'root',
      password: 'password',
      database: 'test'
    });

    connection.connect((err) => {
      if (err) throw err;
      console.log('Connected to MySQL!');
    });

    connection.query('SELECT * FROM users', (err, results) => {
      if (err) throw err;
      console.log(results);
    });
    \end{verbatim}
    
    \item E um exemplo para conexão com o MongoDB:
    \begin{verbatim}
    const mongoose = require('mongoose');

    mongoose.connect('mongodb://localhost:27017/test', {useNewUrlParser: true, useUnifiedTopology: true});

    const db = mongoose.connection;
    db.on('error', console.error.bind(console, 'connection error:'));
    db.once('open', () => {
      console.log('Connected to MongoDB!');
    });
    \end{verbatim}
\end{itemize}

% Comentário: Você pode adicionar qualquer instrução ou dica adicional específica para os leitores, dependendo das necessidades do livro.

\cleardoublepage

\chapter{Tutorial Básico de Linhas de Comando no Linux}

Este apêndice oferece um tutorial básico sobre como utilizar a linha de comando no shell do Linux. Ele abrange desde comandos simples para manipulação de arquivos e diretórios até comandos para acessar a rede e usar ferramentas como o cURL para interagir com APIs.

\section{Comandos Básicos de Arquivos e Diretórios}

\subsection{Criando uma Pasta}
Para criar uma nova pasta (diretório), utilize o comando \texttt{mkdir}:
\begin{verbatim}
mkdir nome_da_pasta
\end{verbatim}
Este comando criará um diretório chamado \texttt{nome\_da\_pasta} no local atual.

\subsection{Navegando até uma Pasta}
Para navegar até um diretório específico, use o comando \texttt{cd} (change directory):
\begin{verbatim}
cd nome_da_pasta
\end{verbatim}
Para voltar ao diretório anterior, você pode usar:
\begin{verbatim}
cd ..
\end{verbatim}

\subsection{Criando um Arquivo HTML}
Para criar um novo arquivo HTML, você pode utilizar o comando \texttt{touch}:
\begin{verbatim}
touch index.html
\end{verbatim}
Este comando criará um arquivo vazio chamado \texttt{index.html} no diretório atual.

\subsection{Listando Arquivos e Diretórios}
Para listar todos os arquivos e diretórios no diretório atual, use o comando \texttt{ls}:
\begin{verbatim}
ls
\end{verbatim}
Você pode adicionar opções ao comando \texttt{ls} para listar arquivos com mais detalhes:
\begin{verbatim}
ls -l      # Lista em formato longo, incluindo permissões e tamanhos.
ls -a      # Lista todos os arquivos, incluindo os ocultos (começam com .).
ls -lh     # Lista com tamanhos de arquivo em um formato legível (KB, MB).
\end{verbatim}

\subsection{Procurando Arquivos}
Para procurar arquivos no sistema, use o comando \texttt{find}:
\begin{verbatim}
find /caminho/para/diretorio -name "nome_do_arquivo"
\end{verbatim}
Este comando pesquisa recursivamente a partir do diretório especificado por arquivos que correspondam ao nome fornecido.

\subsection{Verificando o Histórico de Comandos}
Para ver o histórico dos comandos que você executou, utilize o comando \texttt{history}:
\begin{verbatim}
history
\end{verbatim}
Este comando exibe uma lista de todos os comandos que você executou na sessão atual do terminal.

\section{Comandos de Rede e Acesso a Sites Web}

\subsection{Verificando a Configuração de Rede}
Para verificar a configuração de rede da sua máquina, use o comando \texttt{ifconfig} (ou \texttt{ip} nas versões mais recentes do Ubuntu):
\begin{verbatim}
ifconfig        # Exibe as interfaces de rede e suas configurações.
ip addr show    # Comando alternativo mais moderno para mostrar interfaces.
\end{verbatim}

\subsection{Acessando Sites Web}
Para acessar e recuperar dados de um site a partir da linha de comando, você pode usar o comando \texttt{wget}:
\begin{verbatim}
wget https://www.exemplo.com
\end{verbatim}
Este comando baixa o conteúdo da URL especificada para o diretório atual.

\section{Usando o cURL para Acessar APIs}

\subsection{Introdução ao cURL}
O cURL é uma ferramenta de linha de comando que permite realizar requisições a URLs, sendo particularmente útil para acessar APIs RESTful.

\subsection{Realizando uma Requisição GET}
Para realizar uma requisição GET simples a uma API, use o seguinte comando:
\begin{verbatim}
curl https://api.exemplo.com/dados
\end{verbatim}
Este comando envia uma requisição GET para a URL especificada e exibe a resposta no terminal.

\subsection{Enviando Dados com uma Requisição POST}
Para enviar dados em uma requisição POST, você pode usar o cURL da seguinte forma:
\begin{verbatim}
curl -X POST https://api.exemplo.com/usuario \
     -H "Content-Type: application/json" \
     -d '{"nome": "João", "idade": 30}'
\end{verbatim}
Neste exemplo, o comando envia um objeto JSON contendo os dados do usuário para a API.

\subsection{Autenticação com cURL}
Se a API requer autenticação, você pode adicionar um cabeçalho de autorização à sua requisição:
\begin{verbatim}
curl -H "Authorization: Bearer seu_token_aqui" \
     https://api.exemplo.com/dados_protegidos
\end{verbatim}
Este comando envia uma requisição GET autenticada com um token de acesso.

\section{Outros Comandos Úteis}
\begin{itemize}
    \item \textbf{pwd:} Mostra o diretório atual (o caminho completo até onde você está).
    \begin{verbatim}
    pwd
    \end{verbatim}
    
    \item \textbf{rm:} Remove (exclui) arquivos ou diretórios.
    \begin{verbatim}
    rm nome_do_arquivo
    rm -r nome_do_diretorio  # Remove um diretório e todo o seu conteúdo.
    \end{verbatim}
    
    \item \textbf{cp:} Copia arquivos ou diretórios.
    \begin{verbatim}
    cp arquivo_origem arquivo_destino
    cp -r diretorio_origem diretorio_destino  # Copia diretórios recursivamente.
    \end{verbatim}
    
    \item \textbf{mv:} Move ou renomeia arquivos ou diretórios.
    \begin{verbatim}
    mv arquivo_origem arquivo_destino
    mv nome_antigo nome_novo  # Renomeia um arquivo ou diretório.
    \end{verbatim}
    
    \item \textbf{man:} Exibe o manual de um comando, útil para aprender sobre suas opções e uso.
    \begin{verbatim}
    man comando
    \end{verbatim}
\end{itemize}

% Comentário: Este tutorial fornece uma introdução prática para alunos que precisam utilizar o shell do Linux em um ambiente de desenvolvimento web. Comandos adicionais ou mais avançados podem ser incluídos conforme necessário.

\cleardoublepage

% Formatando como em "APÊNDICE A -- Título do apêndice":
% Ref.: https://tex.stackexchange.com/a/479364/91816

\makeatletter
\def\@chapter[#1]#2{\ifnum \c@secnumdepth >\m@ne
    \refstepcounter{chapter}%
    \typeout{\thechapter.}%
    \addcontentsline{toc}{chapter}%
    {\thechapter\space\textendash\space\ #1}% <-- modification
  \else
    \addcontentsline{toc}{chapter}{#1}%
  \fi
  \chaptermark{#1}%
  \addtocontents{lof}{\protect\addvspace{10\p@}}%
  \addtocontents{lot}{\protect\addvspace{10\p@}}%
  \if@twocolumn
    \@topnewpage[\@makechapterhead{#2}]%
  \else
    \@makechapterhead{#2}%
    \@afterheading
  \fi}
\makeatother

% Reiniciando contador de capítulo e seção:

\setcounter{chapter}{0}
\setcounter{section}{0}

% \chapter{Citações}

\thispagestyle{empty} 

Este modelo usa BibTeX para configurar as referências. O arquivo main.bib contém várias entradas de bibliografia como modelos \cite{article,book,booklet,inbook}. Esses modelos podem ser utilizados para incluir outras entradas e citá-las por meio do seguinte comando:
\begin{verbatim}
\cite{nome_da_entrada}
\end{verbatim}

Por exemplo , a entrada
\verbatimfont{\small}
\begin{verbatim}
@article{greenwade93,
    author  = "George D. Greenwade",
    title   = "The {C}omprehensive {T}ex {A}rchive {N}etwork ({CTAN})",
    year    = "1993",
    journal = "TUGBoat",
    volume  = "14",
    number  = "3",
    pages   = "342--351"
}
\end{verbatim}
\verbatimfont{\normalfont}
pode ser citada no texto com
\begin{verbatim}
\cite{greenwade93}
\end{verbatim}
e a citação apareceria assim: \cite{greenwade93}.

Para fazer uma citação direta no formato ABNT, criamos o ambiente \verb|citacao|, que é uma simples generalização do ambiente \verb|quotation| (habilitado por padrão) com um campo específico de autor. Veja o exemplo a seguir:
\begin{verbatim}
\begin{citacao}{Carl Sagan}
    Alegações extraordinárias exigem evidências extraordinárias.
\end{citacao}
\end{verbatim}
Esse código gera uma citação assim:
\begin{citacao}{Carl Sagan}
    Alegações extraordinárias exigem evidências extraordinárias.
\end{citacao}
O comando \verb|\cite{...}| pode ser usado como indicação do autor:
\begin{verbatim}
\begin{citacao}{\cite{greenwade93}}
TEX is a typesetting program designed for high-quality composition of material that contains a lot of mathematical and technical expressions. It has been adopted by many authors and publishers who generate technical books and papers. It was created by Professor Donald E. Knuth of Stanford University, originally for preparation of his book series ``The Art of Computer Programming''. TEX has been made freely available by Knuth.
\end{citacao}
\end{verbatim}
Naturalmente, a referência \verb|grennwade93| deve estar definida no arquivo BibTeX (aqui, \verb|main.bib|). Confira o resultado:
\begin{citacao}{\cite{greenwade93}}
TEX is a typesetting program designed for high-quality composition of material that contains a lot of mathematical and technical expressions. It has been adopted by many authors and publishers who generate technical books and papers. It was created by Professor Donald E. Knuth of Stanford University, originally for preparation of his book series ``The Art of Computer Programming''. TEX has been made freely available by Knuth.
\end{citacao}



% \cleardoublepage

% \chapter{Tabelas}

\thispagestyle{empty} 

A tabela \ref{tab:SI-basicas} mostra algumas unidades básicas do SI.

\begin{table}
\begin{minipage}{\hsize}
\begin{center}
	\caption{Algumas unidades básicas do SI.}
	\label{tab:SI-basicas}
	\begin{tabular}{P{0.40\hsize}|P{0.5\hsize}}
		\hline
		\textbf{Grandeza} & \textbf{Unidade} \\
		\hline
		comprimento & metro (\si{\meter}) \\
		\hline
		massa & quilograma (\si{\kilo\gram}) \\
		\hline
		tempo & segundo (\si{\second}) \\
		\hline
		corrente elétrica & ampère (\si{\ampere}) \\
		\hline
		temperatura & kelvin (\si{\kelvin}) \\
		\hline
		quantidade de matéria & mol (\si{\mol}) \\
		\hline
		intensidade luminosa & candela (\si{\candela}) \\
		\hline
	\end{tabular}
	\source{adaptado do livro Física para Ciências Agrárias e Ambientais, de Leonardo Luiz e Castro e Olavo Leopoldino da Silva Filho.}
\end{center}
\end{minipage}
\end{table}


É uma boa ideia usar o pacote ``longtable'' para criar tabelas, \index{tabelas} pois assim uma mesma tabela pode ocupar várias páginas. Dê uma olhada no código da lista de símbolos, pois ela foi feita com esse pacote.

% \cleardoublepage

\end{appendices}

% Anexos:

% \renewcommand\appendixname{ANEXO}
% \renewcommand\appendixpagename{ANEXOS}
% \renewcommand{\appendixtocname}{ANEXOS}

% \begin{appendices}

% % Formatando como em "APÊNDICE A -- Título do apêndice":
% Ref.: https://tex.stackexchange.com/a/479364/91816

\makeatletter
\def\@chapter[#1]#2{\ifnum \c@secnumdepth >\m@ne
    \refstepcounter{chapter}%
    \typeout{\thechapter.}%
    \addcontentsline{toc}{chapter}%
    {\thechapter\space\textendash\space\ #1}% <-- modification
  \else
    \addcontentsline{toc}{chapter}{#1}%
  \fi
  \chaptermark{#1}%
  \addtocontents{lof}{\protect\addvspace{10\p@}}%
  \addtocontents{lot}{\protect\addvspace{10\p@}}%
  \if@twocolumn
    \@topnewpage[\@makechapterhead{#2}]%
  \else
    \@makechapterhead{#2}%
    \@afterheading
  \fi}
\makeatother

% Reiniciando contador de capítulo e seção:

\setcounter{chapter}{0}
\setcounter{section}{0}

% \chapter{Figuras de exemplo}

\thispagestyle{empty} 

% Fonte: https://tex.stackexchange.com/a/231741/91816 (Gonzalo Medina)

\noindent\includegraphics[width=3cm]{example-image-a}\qquad
\includegraphics[width=3cm]{example-image-golden}\qquad
\includegraphics[width=3cm]{example-grid-100x100pt}

\noindent\includegraphics[height=5cm]{example-image-b} 

\noindent\includegraphics[scale=0.5]{example-image-c} 

\noindent\includegraphics[width=3cm]{example-image} 
% \cleardoublepage

% \chapter{Logo LaTeX}

\thispagestyle{empty} 

\vfill

\begin{center}
{\fontsize{50pt}{800pt}\selectfont \LaTeX}\
\end{center}

\vfill

\begin{center}
{\color{cinza_UnB}\fontsize{50pt}{800pt}\selectfont \LaTeX}\
\end{center}

\vfill

\begin{center}
{\color{verde_UnB}\fontsize{50pt}{800pt}\selectfont \LaTeX}\
\end{center}

\vfill
% \cleardoublepage

% \end{appendices}

% Parte Pós-Textual:
\backmatter

% Referências (Bibliografia):

\singlespacing

\bibliographystyle{abntex2-alf}
\bibliography{main}
\cleardoublepage

% Índice Remissivo, controlado pelo comando \index{...} em meio ao texto:

\phantomsection
\printindex

\end{document}
